\documentclass{book}
\usepackage{amsmath,amssymb}

\begin{document}

\tableofcontents

\chapter{Introduzione alle teniche numeriche usate nei modelli dinamici}

\section{Una tassonomia dei modelli dinamici}

Un modello dinamico descrive l'evoluzione nel tempo di una o più variabili che rappresentano lo stato di un sistema economico o finanziario. A differenza di un modello statico, nel quale l'attenzione è concentrata sulle relazioni tra variabili in un determinato istante, un modello dinamico permette di studiare una \emph{traiettoria}, ossia una successione di stati del sistema nel corso del tempo.

La varietà dei modelli dinamici utilizzati in economia e finanza è molto ampia. Prima di affrontare i metodi computazionali necessari per risolverli è quindi utile introdurre una tassonomia che permetta di distinguere le principali classi di problemi. La classificazione adottata in questo capitolo si basa su alcune caratteristiche fondamentali: la natura discreta o continua del tempo, la presenza o meno di variabili di controllo, la natura discreta o continua delle variabili di stato e di controllo e, infine, la presenza o meno di incertezza.

Queste caratteristiche non rappresentano semplicemente modi diversi di scrivere uno stesso modello. Esse determinano la struttura matematica del problema e, di conseguenza, i metodi computazionali più appropriati per risolverlo.

\subsection{Modelli dinamici senza controllo}

La classe più semplice è costituita dai modelli nei quali la legge di evoluzione del sistema è assegnata. Una volta conosciuto lo stato iniziale, non vi è alcuna decisione da prendere: l'obiettivo è determinare la traiettoria generata dalla legge di movimento.

La prima distinzione riguarda la rappresentazione del tempo.

In un modello a \emph{tempo discreto}, il tempo assume valori appartenenti a una successione numerabile,

$$
t=0,1,2,\ldots
$$

e una semplice legge di movimento può essere rappresentata nella forma

$$
x_{t+1}=F(x_t).
$$

Conoscendo $x_0$, il modello determina successivamente

$$
x_1,\;x_2,\;x_3,\ldots
$$

La soluzione può quindi essere ottenuta, nei casi più semplici, mediante iterazione in avanti. Modelli di questo tipo sono particolarmente naturali quando il fenomeno considerato presenta una periodicità economicamente significativa, come nel caso di decisioni annuali di investimento, cicli produttivi, pagamenti periodici o osservazioni associate a periodi ben definiti.

In un modello a \emph{tempo continuo}, invece, il tempo può assumere qualsiasi valore reale,

$$
t\in\mathbb{R}_{+},
$$

e l'evoluzione dello stato viene generalmente descritta mediante un'equazione differenziale,

$$
\dot{x}(t)=f(x(t)).
$$

Assegnata la condizione iniziale

$$
x(0)=x_0,
$$

il problema consiste nel determinare la funzione $x(t)$ che descrive l'intera traiettoria del sistema. Questa formulazione è naturale quando si vuole rappresentare un processo che può modificarsi in ogni istante, come l'accumulazione di uno stock, la dinamica di una risorsa o l'evoluzione di alcune variabili finanziarie.

La scelta tra tempo discreto e continuo è quindi innanzitutto una scelta di modellizzazione. Non esiste una formulazione universalmente preferibile. Il tempo discreto può essere più naturale quando la periodicità delle decisioni costituisce una caratteristica essenziale del problema; il tempo continuo può risultare più appropriato quando interessa rappresentare tassi istantanei di variazione o quando una suddivisione del tempo in periodi sarebbe prevalentemente artificiale.

È inoltre importante distinguere la frequenza con cui un fenomeno viene osservato dalla natura temporale del modello. Il fatto che una variabile sia disponibile, ad esempio, con frequenza trimestrale non implica necessariamente che il processo economico sottostante si modifichi soltanto una volta ogni trimestre. La frequenza dei dati e la rappresentazione temporale del modello sono pertanto concetti distinti.

La scelta ha anche conseguenze computazionali. Una legge di movimento discreta ed esplicita può essere direttamente iterata. Un'equazione differenziale deve invece essere generalmente approssimata numericamente su una griglia temporale. Ad esempio, applicando il metodo di Euler a

$$
\dot{x}=f(x)
$$

si ottiene

$$
x_{n+1}=x_n+h f(x_n),
$$

dove $h$ rappresenta il passo temporale utilizzato dall'algoritmo. Il modello economico rimane formulato in tempo continuo, mentre la sua soluzione numerica viene calcolata su un insieme discreto di punti. È quindi necessario distinguere il \emph{tempo del modello} dalla discretizzazione utilizzata dal metodo computazionale.

\subsection{Modelli dinamici con controllo}

Una seconda classe di problemi emerge quando l'evoluzione del sistema non è completamente determinata dalla legge di movimento, ma dipende anche dalle decisioni di uno o più agenti.

In tempo discreto possiamo scrivere

$$
x_{t+1}=F(x_t,u_t),
$$

mentre in tempo continuo possiamo avere

$$
\dot{x}(t)=f(x(t),u(t)).
$$

La variabile $x$ è detta \emph{variabile di stato}, mentre $u$ è detta \emph{variabile di controllo}. Lo stato descrive la situazione nella quale si trova il sistema, mentre il controllo rappresenta una decisione che può modificarne l'evoluzione.

La presenza di una variabile di controllo modifica profondamente la natura del problema. Nei modelli senza controllo la domanda fondamentale era

\begin{quote}
Data la legge di movimento, quale traiettoria seguirà il sistema?
\end{quote}

Nei modelli con controllo la domanda diventa invece

\begin{quote}
Quali decisioni devono essere prese affinché il sistema segua la traiettoria desiderata o ottimale?
\end{quote}

Alla distinzione tra tempo discreto e continuo si aggiunge quindi quella relativa alla natura della variabile di controllo.

Un controllo è \emph{discreto} quando l'agente può scegliere tra un insieme discreto di alternative. Ad esempio,

$$
u_t\in\{0,1\}
$$

può rappresentare la decisione di investire o non investire, entrare o non entrare in un mercato, sostituire o non sostituire un macchinario oppure esercitare o non esercitare un'opzione.

Un controllo è invece \emph{continuo} quando può assumere valori appartenenti a un intervallo,

$$
u_t\in U\subseteq\mathbb{R}^m.
$$

Il consumo, l'investimento, la quantità prodotta o la quota di ricchezza investita in una determinata attività finanziaria costituiscono esempi tipici di controlli continui.

È importante non confondere la natura del tempo con quella del controllo. Un modello può essere formulato in tempo discreto e avere un controllo continuo, oppure essere formulato in tempo continuo e prevedere decisioni discrete. Le due classificazioni descrivono caratteristiche differenti del problema.

\subsection{La natura della variabile di stato}

Nei modelli con controllo diventa inoltre particolarmente rilevante distinguere tra stati discreti e continui.

Una variabile di stato discreta assume valori appartenenti a un insieme finito o numerabile,

$$
x_t\in\{x_1,x_2,\ldots,x_N\}.
$$

Essa può rappresentare, ad esempio, lo stato occupazionale di un individuo, il rating di un'obbligazione, lo stato operativo di un macchinario o un regime economico.

Una variabile di stato continua assume invece valori in un sottoinsieme di $\mathbb{R}^n$,

$$
x_t\in X\subseteq\mathbb{R}^n,
$$

come avviene nel caso della ricchezza, del capitale, del debito o dello stock di una risorsa.

Questa distinzione è particolarmente importante dal punto di vista computazionale. Se lo spazio degli stati è finito, una funzione che associa un valore a ciascuno stato può essere rappresentata mediante un vettore. Se lo stato è continuo, occorre invece rappresentare numericamente un'intera funzione. Diventano allora importanti tecniche quali l'interpolazione, l'approssimazione di funzioni e i metodi di collocazione.

La natura dello stato e quella del controllo costituiscono pertanto ulteriori dimensioni della classificazione dei modelli dinamici e contribuiscono a determinare il metodo computazionale più appropriato.

\subsection{L'introduzione dell'incertezza}

Finora abbiamo implicitamente considerato modelli deterministici. In un modello deterministico, una volta specificati lo stato iniziale e, quando presenti, i controlli, l'evoluzione futura del sistema è completamente determinata.

L'incertezza può essere introdotta come un ulteriore livello di complessità che si sovrappone alla classificazione precedente. Non costituisce quindi una categoria alternativa ai modelli con o senza controllo: entrambe le classi possono essere formulate in versione deterministica oppure stocastica.

Ad esempio, una legge di movimento discreta senza controllo,

$$
x_{t+1}=F(x_t),
$$

può essere estesa introducendo uno shock casuale,

$$
x_{t+1}=F(x_t,\varepsilon_{t+1}),
$$

dove $\varepsilon_{t+1}$ è una variabile aleatoria.

Analogamente, in tempo continuo, una dinamica deterministica della forma

$$
dx_t=\mu(x_t)\,dt
$$

può essere trasformata in un processo stocastico introducendo un termine di diffusione,

$$
dx_t=\mu(x_t)\,dt+\sigma(x_t)\,dW_t,
$$

dove $W_t$ rappresenta un moto browniano.

L'incertezza può naturalmente essere introdotta anche nei modelli con controllo. In tempo discreto possiamo avere

$$
x_{t+1}=F(x_t,u_t,\varepsilon_{t+1}),
$$

mentre in tempo continuo una possibile formulazione è

$$
dx_t=
\mu(x_t,u_t)\,dt+
\sigma(x_t,u_t)\,dW_t.
$$

In questi casi l'agente deve prendere decisioni senza conoscere con certezza gli stati futuri del sistema. Il problema computazionale non riguarda quindi più soltanto l'evoluzione dinamica e l'ottimizzazione, ma anche la rappresentazione e l'integrazione rispetto all'incertezza.

\subsection{Una classificazione multidimensionale}

La tassonomia appena descritta può essere interpretata come una classificazione costruita lungo dimensioni differenti e, almeno in parte, indipendenti.

Una prima dimensione riguarda la presenza di decisioni:

$$
\text{senza controllo}
\qquad\text{oppure}\qquad
\text{con controllo}.
$$

Una seconda dimensione riguarda la rappresentazione del tempo:

$$
\text{tempo discreto}
\qquad\text{oppure}\qquad
\text{tempo continuo}.
$$

Nei modelli con controllo diventano inoltre rilevanti la natura dello stato,

$$
\text{stato discreto}
\qquad\text{oppure}\qquad
\text{stato continuo},
$$

e la natura del controllo,

$$
\text{controllo discreto}
\qquad\text{oppure}\qquad
\text{controllo continuo}.
$$

Infine, a ciascuna delle precedenti configurazioni può essere aggiunta l'incertezza:

$$
\text{deterministico}
\qquad\longrightarrow\qquad
\text{stocastico}.
$$

Non tutte le possibili combinazioni verranno studiate con lo stesso livello di dettaglio. Lo scopo della tassonomia è piuttosto fornire una mappa concettuale che permetta di collocare i diversi modelli e, soprattutto, di comprendere perché essi richiedano tecniche computazionali differenti.

La progressione seguita nei capitoli successivi sarà pertanto graduale. Si partirà dai modelli deterministici senza controllo, studiando separatamente il tempo discreto e quello continuo. Si introdurranno successivamente le decisioni e i problemi di controllo dinamico. Infine, l'incertezza verrà considerata come un ulteriore elemento che può essere incorporato sia nei modelli dinamici senza controllo sia nei problemi di decisione dinamica.

Questa progressione può essere sintetizzata come

$$
\boxed{
\text{dinamica}
\;\longrightarrow\;
\text{decisione}
\;\longrightarrow\;
\text{incertezza}
}
$$

e costituisce anche una progressione dal punto di vista computazionale. A ogni passaggio, infatti, nuovi problemi numerici si aggiungono a quelli già incontrati: l'iterazione e l'integrazione numerica per determinare le traiettorie, l'ottimizzazione e la soluzione di problemi con condizioni al contorno quando vengono introdotte le decisioni, e infine la simulazione, l'integrazione rispetto a distribuzioni di probabilità e l'approssimazione di funzioni quando viene introdotta l'incertezza.

La classificazione dei modelli dinamici non costituisce quindi soltanto un'organizzazione formale dei diversi casi. Essa permette di stabilire un collegamento diretto tra la struttura economica del modello e gli strumenti computazionali necessari per studiarlo.




\section{Tempo discreto e tempo continuo}

La tassonomia introdotta nella sezione precedente distingue i modelli dinamici secondo diverse dimensioni. Alcune di queste caratteristiche sono in larga misura determinate dalla natura economica del problema. Altre dipendono maggiormente dalle scelte effettuate dal modellista. La distinzione tra tempo discreto e tempo continuo appartiene soprattutto a questa seconda categoria.

Si consideri, ad esempio, la natura di una variabile di controllo. Se un'impresa deve decidere se entrare oppure non entrare in un mercato, il controllo ha una natura intrinsecamente discreta:

$$
u\in\{0,1\}.
$$

Se invece l'impresa deve scegliere la quantità da produrre, la decisione può essere naturalmente rappresentata mediante una variabile continua:

$$
q\in\mathbb{R}_{+}.
$$

Considerazioni analoghe valgono per le variabili di stato. Il rating di un'obbligazione appartiene a un insieme discreto di categorie, mentre la ricchezza di un investitore, il capitale di un'impresa o lo stock di una risorsa vengono generalmente rappresentati mediante variabili continue.

Naturalmente anche queste rappresentazioni contengono una componente di modellizzazione. Una variabile teoricamente continua può essere discretizzata e una realtà complessa può essere rappresentata mediante un numero finito di stati. Tuttavia, la natura economica della variabile fornisce spesso un'indicazione abbastanza chiara sulla rappresentazione più appropriata.

La scelta tra un modello a tempo discreto e uno a tempo continuo è invece generalmente meno immediata. Uno stesso fenomeno economico può spesso essere rappresentato in entrambi i modi. L'accumulazione del capitale, la dinamica della ricchezza, l'evoluzione del debito, la crescita economica e la dinamica dei prezzi finanziari forniscono esempi nei quali entrambe le formulazioni sono possibili.

La domanda rilevante non è quindi quale delle due rappresentazioni sia in assoluto corretta, ma quale sia più appropriata rispetto al fenomeno studiato, alla domanda di ricerca e agli strumenti analitici e computazionali che si intendono utilizzare.

\subsection{Due rappresentazioni del tempo}

In un modello a tempo discreto il tempo assume valori appartenenti a un insieme numerabile:

$$
t=0,1,2,\ldots
$$

e la dinamica può essere descritta, nella forma più semplice, mediante

$$
x_{t+1}=F(x_t).
$$

Il passaggio da $t$ a $t+1$ rappresenta un intervallo temporale la cui durata deve essere definita dal modellista: un giorno, un mese, un trimestre, un anno o qualsiasi altra unità appropriata.

In un modello a tempo continuo, invece,

$$
t\in\mathbb{R}_{+},
$$

e la dinamica viene tipicamente rappresentata attraverso il tasso istantaneo di variazione dello stato:

$$
\dot{x}(t)=f(x(t)).
$$

In questo caso non viene imposto un intervallo fondamentale tra due date consecutive. Il sistema viene idealmente considerato in grado di modificarsi in ogni istante.

Questa differenza apparentemente formale riflette due diverse rappresentazioni del processo economico. Nel primo caso l'evoluzione viene descritta attraverso transizioni tra periodi; nel secondo attraverso variazioni istantanee.

Nella scelta del tipo di modello da usare, pesano in modo particolare considerazioni relative alla natura economica del fenomeno e alla frequenza temporale e dei dati con le quali può essere rappresentato.

\subsubsection{La natura economica del fenomeno}

Un primo criterio per scegliere tra le due formulazioni riguarda il modo in cui il fenomeno economico evolve nel tempo.

Il tempo discreto è particolarmente naturale quando esiste una periodicità economicamente significativa. Si pensi a una decisione di semina effettuata una volta per stagione, alla sostituzione annuale di un macchinario, al pagamento di una cedola in date prestabilite o a una decisione di politica economica presa periodicamente. In questi casi il periodo può corrispondere a una caratteristica effettiva del processo decisionale.

Il tempo continuo può risultare più naturale quando non esiste una frequenza fondamentale degli aggiustamenti. Il valore di un'attività finanziaria, uno stock di capitale, una risorsa naturale o un saldo monetario possono, almeno idealmente, modificarsi in qualsiasi momento. Una rappresentazione continua evita allora di introdurre artificialmente una periodicità che non appartiene necessariamente al fenomeno.

Tuttavia, questa distinzione non deve essere interpretata in maniera troppo rigida. Molte decisioni che nella realtà avvengono in date specifiche possono essere convenientemente approssimate mediante modelli continui; analogamente, processi potenzialmente continui possono essere rappresentati mediante modelli discreti.

La scelta della rappresentazione temporale costituisce quindi un compromesso tra realismo, semplicità e finalità del modello.

\subsubsection{Frequenza del modello e frequenza dei dati}

Una distinzione particolarmente importante riguarda la frequenza delle osservazioni disponibili.

Il fatto che una variabile sia osservata trimestralmente non implica che il processo economico sottostante evolva soltanto una volta ogni trimestre. Analogamente, la disponibilità di dati giornalieri non obbliga a formulare un modello con periodo giornaliero.

In generale,

$$
\boxed{
\text{frequenza delle osservazioni}
\neq
\text{frequenza del processo economico}
}.
$$

Un processo formulato in tempo continuo viene necessariamente osservato a intervalli discreti. Allo stesso modo, un modello teorico discreto può utilizzare periodi differenti dalla frequenza alla quale i dati sono originariamente raccolti.

La frequenza dei dati rappresenta quindi un elemento da considerare nella scelta del modello, ma non costituisce da sola un criterio sufficiente per determinare la rappresentazione temporale.

\subsection{Il tempo discreto}
\subsubsection{Vantaggi del tempo discreto}

Uno dei principali vantaggi del tempo discreto è la semplicità della struttura ricorsiva. Se la legge di movimento è esplicita,

$$
x_{t+1}=F(x_t),
$$

e lo stato iniziale $x_0$ è noto, la traiettoria può essere costruita mediante successive iterazioni:

$$
x_0
\longrightarrow
x_1
\longrightarrow
x_2
\longrightarrow
\cdots.
$$

In questo caso non è necessaria alcuna discretizzazione numerica del tempo: la struttura discreta appartiene al modello stesso.

Il tempo discreto rende inoltre particolarmente trasparente il timing degli eventi. È possibile specificare se una decisione viene presa prima o dopo l'osservazione di uno shock, se gli interessi vengono pagati prima o dopo un prelievo, oppure se una decisione di investimento influenza la capacità produttiva corrente o quella del periodo successivo.

Questa caratteristica rende la formulazione discreta particolarmente naturale per molti modelli ricorsivi e per il dynamic programming.

Un ulteriore vantaggio è la vicinanza con numerosi dati economici e finanziari, che vengono osservati a frequenze discrete. Sebbene, come osservato in precedenza, ciò non imponga una formulazione discreta, può facilitare l'interpretazione, la calibrazione e la stima del modello.

\subsubsection{Limiti del tempo discreto}

La formulazione discreta richiede tuttavia di scegliere la durata del periodo. Tale scelta può essere economicamente rilevante.

Si consideri, ad esempio, una semplice equazione di accumulazione:

$$
K_{t+1}=(1-\delta)K_t+I_t.
$$

Il parametro $\delta$ rappresenta il tasso di deprezzamento per periodo. Il suo significato dipende quindi dalla durata del periodo adottato. Un tasso annuale non può essere utilizzato senza modifiche in un modello mensile.

Più in generale, modificare la durata del periodo può richiedere una trasformazione dei parametri e, in alcuni casi, può modificare anche la struttura economica del modello.

Il tempo discreto richiede inoltre di specificare l'ordine degli eventi all'interno del periodo. Ad esempio,

$$
a_{t+1}=(1+r)a_t+y_t-c_t
$$

e

$$
a_{t+1}=(1+r)(a_t+y_t-c_t)
$$

non descrivono lo stesso timing economico. Nel primo caso soltanto la ricchezza iniziale produce il rendimento $r$ durante il periodo; nel secondo anche il reddito netto del consumo viene investito allo stesso rendimento.

Quando la suddivisione del tempo è puramente artificiale, tali scelte di timing possono diventare esse stesse artificiose.


\subsection{Il tempo continuo}
\subsubsection{Vantaggi del tempo continuo}

La formulazione continua descrive direttamente il tasso istantaneo di variazione di una variabile. Ad esempio,

$$
\dot K(t)=I(t)-\delta K(t)
$$

separa immediatamente il flusso di investimento dal deprezzamento istantaneo dello stock di capitale.

Il tempo continuo permette inoltre di utilizzare gli strumenti dell'analisi differenziale e della teoria delle equazioni differenziali. Concetti quali stabilità, sistemi dinamici, controllo ottimo e, in presenza di incertezza, equazioni differenziali stocastiche possono essere formulati in maniera particolarmente naturale.

In alcuni problemi la formulazione continua può anche migliorare la trattabilità matematica. Relazioni che in tempo discreto coinvolgono condizioni tra periodi successivi possono diventare condizioni locali basate su derivate. Ciò spiega in parte l'ampio utilizzo del tempo continuo in macroeconomia, finanza e teoria del controllo.

In finanza, in particolare, il continuous time fornisce un linguaggio naturale per rappresentare prezzi che possono modificarsi in ogni istante e per sviluppare argomenti di arbitraggio, hedging e valutazione dei derivati.

\subsubsection{Limiti del tempo continuo}

La maggiore eleganza analitica non implica necessariamente una maggiore semplicità computazionale.

Un computer non può rappresentare direttamente un continuum di istanti. Per calcolare numericamente la soluzione di

$$
\dot x(t)=f(x(t)),
$$

è necessario considerare una griglia

$$
t_0,t_1,\ldots,t_N
$$

e approssimare la soluzione nei punti della griglia.

Ad esempio, il metodo di Euler produce

$$
x_{n+1}=x_n+h f(x_n),
$$

dove $h$ rappresenta il passo utilizzato per la discretizzazione.

Questo introduce una distinzione fondamentale:

$$
\boxed{
\text{modello continuo}
\neq
\text{algoritmo continuo}
}.
$$

Il modello può essere formulato in tempo continuo, ma la sua soluzione numerica viene necessariamente approssimata utilizzando un numero finito di operazioni.

Di conseguenza, nella soluzione numerica dei modelli continui compare un problema che non esiste nella stessa forma quando si itera una legge di movimento originariamente discreta: l'errore di discretizzazione temporale. La scelta del passo $h$ implica generalmente un compromesso tra accuratezza e costo computazionale.

È tuttavia importante non concludere che i modelli continui siano necessariamente più difficili da risolvere. In alcune applicazioni la formulazione continua produce equazioni con una struttura matematica particolarmente favorevole e, dopo la discretizzazione, sistemi sparsi che possono essere risolti in modo molto efficiente. La convenienza computazionale deve quindi essere valutata con riferimento allo specifico problema.

\subsection{Il legame tra le due rappresentazioni}

Tempo discreto e tempo continuo non costituiscono due mondi completamente separati. In molti casi esiste una relazione molto stretta tra le due formulazioni.

Consideriamo

$$
\dot x(t)=f(x(t)).
$$

Per un piccolo intervallo temporale $h$,

$$
\frac{x(t+h)-x(t)}{h}
\approx
f(x(t)),
$$

da cui

$$
x(t+h)
\approx
x(t)+h f(x(t)).
$$

Ponendo

$$
x_n=x(t_n),
$$

otteniamo

$$
x_{n+1}=x_n+h f(x_n).
$$

La discretizzazione produce quindi una legge di movimento in tempo discreto. Al diminuire di $h$, sotto opportune condizioni, la traiettoria discreta può convergere alla soluzione del modello continuo.

Questa relazione è importante, ma richiede una precisazione. Una equazione alle differenze ottenuta discretizzando un'equazione differenziale non deve essere confusa con un modello formulato originariamente in tempo discreto. Nel primo caso il periodo $h$ è un parametro dell'algoritmo numerico; nel secondo la durata del periodo appartiene alla struttura economica del modello.

\subsubsection{Quando discreto e continuo producono risultati simili}

In numerosi problemi le due rappresentazioni possono descrivere sostanzialmente lo stesso meccanismo economico.

Un semplice modello di aggiustamento può, ad esempio, essere scritto in tempo continuo come

$$
\dot x(t)=\alpha\left[x^*-x(t)\right],
$$

oppure, in tempo discreto, come

$$
x_{t+1}-x_t
=
\gamma\left(x^*-x_t\right).
$$

Entrambi descrivono un processo nel quale lo stato si muove verso $x^*$. Scegliendo opportunamente la relazione tra $\alpha$, $\gamma$ e la durata del periodo, le due formulazioni possono generare traiettorie molto simili.

Lo stesso principio emerge in molti modelli di crescita, accumulazione della ricchezza e dinamica finanziaria. Quando il periodo discreto è sufficientemente breve rispetto alla velocità con cui evolve il fenomeno e i parametri vengono trasformati in maniera coerente, il modello discreto può fornire una buona approssimazione della corrispondente dinamica continua.

Questa possibilità spiega perché numerosi meccanismi economici possano essere studiati indifferentemente, almeno per alcuni obiettivi, utilizzando una delle due rappresentazioni.

\subsubsection{Quando discreto e continuo possono produrre risultati diversi}

La possibilità di collegare le due formulazioni non implica tuttavia una loro equivalenza generale.

Innanzitutto, una diversa struttura temporale può modificare il timing delle decisioni e dell'informazione. In un modello discreto può essere essenziale stabilire se uno shock viene osservato prima o dopo una decisione. Nel limite continuo alcune di queste distinzioni assumono una forma differente.

In secondo luogo, le proprietà dinamiche possono dipendere dalla frequenza di aggiustamento. Una dinamica discreta può, ad esempio, generare oscillazioni o comportamenti complessi che non emergono necessariamente dalla corrispondente equazione differenziale.

Inoltre, in modelli con aspettative, vincoli finanziari, non linearità o decisioni forward-looking, il passaggio da una formulazione all'altra può modificare le proprietà quantitative e, in alcuni casi, qualitative dell'equilibrio.

Infine, soprattutto nella finanza stocastica, il tempo continuo consente di costruire processi e strutture matematiche che non possiedono necessariamente un equivalente discreto esatto.

Pertanto,

\begin{quotation}
tempo discreto e tempo continuo possono essere approssimativamente equivalenti in alcuni problemi,
\end{quotation}
ma
\begin{quotation}
non sono rappresentazioni universalmente equivalenti.
\end{quotation}


La scelta della struttura temporale può quindi avere conseguenze sulle conclusioni economiche del modello e non deve essere considerata esclusivamente una questione di convenienza matematica.

\subsection{Considerazioni finali}

Non esiste una regola generale che permetta di stabilire che il tempo discreto sia preferibile al tempo continuo, o viceversa. La scelta dovrebbe derivare dalla valutazione congiunta di diversi elementi.

Il primo riguarda il \emph{meccanismo economico}. Se esiste una periodicità intrinseca nelle decisioni o negli eventi, una formulazione discreta può essere più naturale. Se il processo può modificarsi continuamente e la periodicità sarebbe prevalentemente artificiale, il tempo continuo può essere più appropriato.

Il secondo riguarda la \emph{domanda di ricerca}. Un modello costruito per studiare decisioni annuali di investimento può richiedere una rappresentazione diversa da un modello costruito per valutare un derivato finanziario.

Il terzo riguarda la \emph{trattabilità}. In alcuni problemi il tempo discreto produce una struttura ricorsiva semplice; in altri il passaggio al continuo consente una formulazione matematica più compatta e trattabile.

Il quarto riguarda la \emph{soluzione computazionale}. È necessario considerare quali algoritmi siano disponibili, il grado di accuratezza richiesto, la dimensionalità del problema e il costo della soluzione.

Infine, occorre considerare il collegamento con i \emph{dati}. Sebbene la frequenza delle osservazioni non determini automaticamente la frequenza del modello, la disponibilità e la struttura dei dati possono rendere una formulazione più conveniente dell'altra.

La scelta tra tempo discreto e continuo deve pertanto precedere la scelta dell'algoritmo numerico. Prima si stabilisce quale rappresentazione del tempo sia più appropriata per il fenomeno che si intende studiare; successivamente si individua il metodo computazionale adatto a risolvere il modello così formulato.

Questo principio sarà importante nelle sezioni successive. Inizieremo studiando separatamente i modelli dinamici deterministici a tempo discreto e a tempo continuo. Il confronto tra le due formulazioni permetterà non soltanto di comprendere i diversi metodi numerici, ma anche di valutare come una scelta apparentemente tecnica sulla rappresentazione del tempo possa influenzare la costruzione e le implicazioni di un modello economico o finanziario.


\section{Illustrazione dei contenuti di questa parte}

TBW


\part{Modelli in tempo discreto senza controllo}


Molti problemi economici e finanziari sono formulati in tempo discreto.
Il tempo è rappresentato da una successione di periodi,
\[
t=0,1,2,\ldots,
\]
e l'evoluzione del sistema è descritta attraverso relazioni che collegano
le variabili economiche osservate in periodi differenti.

La formulazione più semplice è quella già introdotta in precedenza:
\[
x_{t+1}=F(x_t),
\]
dove $x_t$ rappresenta lo stato del sistema al tempo $t$ e $F$ ne
descrive la legge di movimento. Una volta assegnato lo stato iniziale
$x_0$, il modello può essere iterato in avanti:
\[
x_0 \longrightarrow x_1 \longrightarrow x_2
\longrightarrow \cdots.
\]

Quando il sistema contiene più variabili, è conveniente raccoglierle
nel vettore
\[
\mathbf{x}_t=
\begin{bmatrix}
x_{1,t}\\
x_{2,t}\\
\vdots\\
x_{n,t}
\end{bmatrix}.
\]
La legge di movimento assume allora la forma
\[
\mathbf{x}_{t+1}=F(\mathbf{x}_t).
\]

Questa rappresentazione, tuttavia, non esaurisce le possibili
formulazioni dei modelli dinamici discreti. Dal punto di vista
computazionale è utile distinguere almeno tre casi: sistemi espliciti
o ricorsivi, sistemi impliciti e sistemi forward-looking.

\paragraph{Sistemi espliciti o ricorsivi}

Nel caso più semplice, le variabili del periodo successivo sono
espresse direttamente in funzione delle variabili già note:
\[
\mathbf{x}_{t+1}=F(\mathbf{x}_t).
\]

Dato $\mathbf{x}_t$, è quindi possibile calcolare direttamente
$\mathbf{x}_{t+1}$. La simulazione consiste semplicemente
nell'applicazione ripetuta della legge di movimento:
\[
\mathbf{x}_0
\rightarrow
\mathbf{x}_1
\rightarrow
\mathbf{x}_2
\rightarrow\cdots.
\]

Nel caso lineare la legge di movimento può essere scritta come
\[
\mathbf{x}_{t+1}
=
A\mathbf{x}_t+\mathbf{b},
\]
dove $A$ è la matrice che descrive le relazioni dinamiche tra le
variabili e $\mathbf{b}$ raccoglie i termini costanti.

Il problema computazionale fondamentale consiste quindi nella
\emph{forward iteration}: a ogni periodo si utilizza lo stato corrente
per calcolare quello successivo.

\paragraph{Sistemi impliciti}

Non tutti i modelli economici forniscono direttamente una legge di
movimento della forma precedente. In molti casi alcune variabili dello
stesso periodo sono determinate simultaneamente.

Una rappresentazione generale è
\[
G(\mathbf{x}_{t+1},\mathbf{x}_t)=0.
\]

Dato $\mathbf{x}_t$, non è quindi possibile ottenere
$\mathbf{x}_{t+1}$ mediante la semplice valutazione di una funzione:
occorre prima risolvere il sistema di equazioni rappresentato da $G$.

Nel caso lineare possiamo avere, ad esempio,
\[
A\mathbf{x}_{t+1}
=
B\mathbf{x}_t+\mathbf{c}.
\]

Se $A$ è non singolare, il sistema ammette la rappresentazione
\[
\mathbf{x}_{t+1}
=
A^{-1}B\mathbf{x}_t+A^{-1}\mathbf{c}.
\]

Dal punto di vista matematico il modello è stato così ricondotto a una
forma ricorsiva. Dal punto di vista computazionale, tuttavia, non è
generalmente conveniente calcolare esplicitamente $A^{-1}$: a ogni
periodo si risolve direttamente il sistema lineare
\[
A\mathbf{x}_{t+1}
=
B\mathbf{x}_t+\mathbf{c}.
\]

Se $G$ è non lineare, la determinazione di $\mathbf{x}_{t+1}$ può
invece richiedere, a ogni periodo, la soluzione numerica di un sistema
di equazioni non lineari.

La distinzione tra sistemi espliciti e impliciti è quindi importante
anche dal punto di vista computazionale:
\[
\begin{array}{ccc}
\text{sistema esplicito}
&
\Rightarrow
&
\text{valutazione della legge di movimento},
\\[2mm]
\text{sistema implicito}
&
\Rightarrow
&
\text{soluzione di un sistema di equazioni}.
\end{array}
\]

\paragraph{Sistemi forward-looking}

Una difficoltà diversa emerge quando le decisioni o le variabili
correnti dipendono da valori futuri del sistema. In questo caso
compaiono variabili \emph{forward-looking}.

In un modello deterministico con previsione perfetta possiamo avere,
ad esempio,
\[
G(\mathbf{x}_{t+1},\mathbf{x}_t,\mathbf{x}_{t-1})=0,
\]
dove alcune condizioni correnti dipendono direttamente da variabili
future.

Quando è presente incertezza, la formulazione tipica coinvolge invece
aspettative condizionate:
\[
G(E_t\mathbf{x}_{t+1},\mathbf{x}_t,\mathbf{x}_{t-1})=0,
\]
dove $E_t\mathbf{x}_{t+1}$ rappresenta il valore atteso delle variabili
future sulla base dell'informazione disponibile al tempo $t$.

Questa struttura è particolarmente importante nei modelli con
aspettative razionali. Le aspettative utilizzate dagli agenti devono
essere coerenti con la dinamica generata dal modello stesso.

In questi casi non è generalmente possibile procedere semplicemente
da
\[
\mathbf{x}_t
\quad\text{a}\quad
\mathbf{x}_{t+1}
\]
come nei sistemi ricorsivi. Il problema deve essere risolto tenendo
conto congiuntamente delle condizioni che collegano presente e futuro.

Dal punto di vista computazionale emerge quindi una distinzione
fondamentale:
\begin{quotation}
	le traiettorie generate dai sistemi esprimibili in forma ricorsiva possono essere generate dalle equazioni del sistema
\end{quotation}
mentre, nei modelli forward-looking,

\begin{quotation}
	occorre risolvere il problema e poi si possono ricavare le traiettorie
\end{quotation}

Occorre cioè determinare prima una soluzione del modello coerente con
le condizioni intertemporali e, successivamente, utilizzare tale
soluzione per generare le traiettorie delle variabili.

\paragraph{Una prima tassonomia computazionale}

Le tre formulazioni possono essere sintetizzate come segue:
\[
\begin{array}{lll}
\textbf{Esplicito:}
&
\mathbf{x}_{t+1}=F(\mathbf{x}_t)
&
\Rightarrow \text{ iterazione diretta},
\\[2mm]
\textbf{Implicito:}
&
G(\mathbf{x}_{t+1},\mathbf{x}_t)=0
&
\Rightarrow \text{ soluzione periodo per periodo},
\\[2mm]
\textbf{Forward-looking:}
&
G(E_t\mathbf{x}_{t+1},\mathbf{x}_t,\ldots)=0
&
\Rightarrow \text{ soluzione intertemporale e successiva simulazione}.
\end{array}
\]

Questa classificazione riguarda il modo in cui la dinamica è
formulata e, di conseguenza, il tipo di problema computazionale che
deve essere risolto. Essa è distinta dalla classificazione tra modelli
deterministici e stocastici.

Un sistema ricorsivo, ad esempio, può essere deterministico,
\[
\mathbf{x}_{t+1}=F(\mathbf{x}_t),
\]
oppure stocastico,
\[
\mathbf{x}_{t+1}
=
F(\mathbf{x}_t,\boldsymbol{\varepsilon}_{t+1}).
\]

Analogamente, un sistema forward-looking può essere studiato in un
ambiente deterministico con previsione perfetta oppure in presenza di
incertezza, nel qual caso le aspettative sul futuro assumono un ruolo
centrale.

Nei prossimi capitoli procederemo quindi gradualmente. Inizieremo dai
modelli discreti deterministici senza controllo, per i quali sarà
possibile osservare direttamente le traiettorie generate dalle
equazioni alle differenze. Vedremo che tali traiettorie possono
cambiare profondamente al variare dei parametri del modello. Questo
condurrà naturalmente al problema della stabilità e, nel caso dei
sistemi multidimensionali, allo studio degli autovalori.

Successivamente introdurremo l'incertezza, passando dai modelli
discreti deterministici ai modelli discreti stocastici e dalle singole
traiettorie alle distribuzioni delle possibili traiettorie.


\paragraph{Condizioni iniziali, condizioni al contorno e traiettorie}

Una volta specificate le equazioni che descrivono la dinamica del sistema,
occorre stabilire quali condizioni siano necessarie per determinare una
particolare traiettoria. La natura di tali condizioni dipende dalla
struttura temporale del modello.

Consideriamo innanzitutto un sistema esplicito della forma
\[
\mathbf{x}_{t+1}=F(\mathbf{x}_t).
\]
Per generare una particolare traiettoria è sufficiente specificare lo
stato iniziale
\[
\mathbf{x}_0.
\]
Il problema costituito dalla legge di movimento e dalla condizione
iniziale prende il nome di \emph{initial value problem} (IVP):
\[
\begin{cases}
\mathbf{x}_{t+1}=F(\mathbf{x}_t),\\
\mathbf{x}_0=\bar{\mathbf{x}}.
\end{cases}
\]

Una volta assegnato $\mathbf{x}_0$, l'iterazione della legge di
movimento genera la traiettoria
\[
\mathbf{x}_0,\mathbf{x}_1,\mathbf{x}_2,\ldots.
\]

La stessa logica si applica ai sistemi formulati in forma implicita:
\[
G(\mathbf{x}_{t+1},\mathbf{x}_t)=0,
\]
purché, dato $\mathbf{x}_t$, le equazioni determinino univocamente
$\mathbf{x}_{t+1}$. Anche in questo caso una condizione iniziale è
sufficiente per determinare la traiettoria. La differenza è di natura
computazionale: mentre nel sistema esplicito lo stato successivo è
ottenuto mediante la valutazione di $F$, nel sistema implicito occorre
risolvere, a ogni periodo, il sistema
\[
G(\mathbf{x}_{t+1},\mathbf{x}_t)=0.
\]

Possiamo quindi rappresentare schematicamente i due casi come
\[
\text{sistema esplicito:}
\qquad
\mathbf{x}_t
\longrightarrow
F
\longrightarrow
\mathbf{x}_{t+1},
\]
e
\[
\text{sistema implicito:}
\qquad
\mathbf{x}_t
\longrightarrow
\text{solver}
\longrightarrow
\mathbf{x}_{t+1}.
\]

In entrambi i casi, una volta noto lo stato iniziale, il problema può
essere risolto procedendo in avanti nel tempo e può quindi essere
interpretato come un IVP.

Una situazione diversa si presenta nei modelli \emph{forward-looking}.
In questi modelli alcune variabili correnti dipendono da condizioni
future. Una rappresentazione schematica può essere
\[
\begin{aligned}
\mathbf{x}_{t+1}
&=
F(\mathbf{x}_t,\mathbf{q}_t),\\
\mathbf{q}_t
&=
G(\mathbf{x}_t,\mathbf{q}_{t+1}),
\end{aligned}
\]
dove $\mathbf{x}_t$ rappresenta un insieme di variabili predeterminate
e $\mathbf{q}_t$ un insieme di variabili forward-looking.

In questo caso conoscere
\[
\mathbf{x}_0
\]
non è generalmente sufficiente per determinare l'intera traiettoria.
Il valore iniziale delle variabili forward-looking non può essere
assegnato arbitrariamente, ma deve essere determinato in modo che la
traiettoria soddisfi opportune condizioni future.

A una condizione iniziale come
\[
\mathbf{x}_0=\bar{\mathbf{x}}
\]
si affianca quindi una condizione terminale, una condizione asintotica,
una condizione di non esplosività oppure, a seconda del modello, una
condizione di trasversalità.

Il problema non è più, in generale, un semplice IVP, ma assume la
struttura di un \emph{boundary value problem} (BVP), nel quale le
condizioni necessarie per selezionare la soluzione non sono concentrate
esclusivamente nel periodo iniziale.

In un problema a orizzonte finito possiamo avere, ad esempio,
\[
\mathbf{x}_0=\bar{\mathbf{x}},
\qquad
\mathbf{q}_T=\bar{\mathbf{q}}.
\]
Alcune condizioni sono quindi specificate all'inizio dell'orizzonte e
altre alla fine. Si parla, in questo caso, di condizioni al contorno
miste.

Nei modelli a orizzonte infinito la condizione futura non deve
necessariamente specificare il valore assunto da una variabile in un
periodo terminale. Può invece imporre un comportamento asintotico della
traiettoria, come una condizione di non esplosività.

La distinzione tra IVP e BVP introduce quindi una prima importante
differenza computazionale:
\[
\boxed{
\begin{array}{rcl}
\text{IVP}
&:&
\text{condizione iniziale}
\longrightarrow
\text{simulazione in avanti},\\[2mm]
\text{BVP}
&:&
\text{condizioni iniziali e future}
\longrightarrow
\text{soluzione della traiettoria}.
\end{array}
}
\]

Nel caso forward-looking occorre quindi distinguere tra la
\emph{soluzione del modello} e la successiva \emph{simulazione della
soluzione}. Una volta risolto il BVP, è spesso possibile ottenere una
regola di soluzione ricorsiva della forma
\[
\mathbf{q}_t=P\mathbf{x}_t
\]
e una legge di movimento
\[
\mathbf{x}_{t+1}=H(\mathbf{x}_t).
\]
La soluzione così ottenuta può nuovamente essere simulata procedendo in
avanti a partire dalla condizione iniziale.

Questa distinzione può essere sintetizzata come
\[
\text{modelli ricorsivi}
\quad\longrightarrow\quad
\text{simulate},
\]
mentre
\[
\text{modelli forward-looking}
\quad\longrightarrow\quad
\text{solve first, then simulate}.
\]


\paragraph{Condizioni iniziali, traiettorie e stabilità}

Nei sistemi che possono essere simulati come IVP, condizioni iniziali
differenti generano, in generale, traiettorie differenti. Se, ad esempio,
\[
\mathbf{x}_0^{(1)}\neq\mathbf{x}_0^{(2)},
\]
le due condizioni iniziali producono
\[
\mathbf{x}_0^{(1)}
\rightarrow
\mathbf{x}_1^{(1)}
\rightarrow
\mathbf{x}_2^{(1)}
\rightarrow\cdots
\]
e
\[
\mathbf{x}_0^{(2)}
\rightarrow
\mathbf{x}_1^{(2)}
\rightarrow
\mathbf{x}_2^{(2)}
\rightarrow\cdots.
\]

La simulazione numerica consente di costruire e visualizzare queste
traiettorie. Tuttavia, l'osservazione di una singola traiettoria non è
sufficiente per comprendere il comportamento generale del sistema.
Occorre chiedersi che cosa accade modificando le condizioni iniziali
e i parametri della legge di movimento.

Due traiettorie differenti possono infatti presentare lo stesso
comportamento qualitativo. Possono, ad esempio, convergere entrambe
verso lo stesso equilibrio, pur partendo da condizioni iniziali molto
diverse. In altri casi, piccole variazioni delle condizioni iniziali o
dei parametri possono modificare profondamente l'evoluzione del
sistema.

Questo conduce a uno dei problemi fondamentali nello studio dei
modelli dinamici: l'\emph{analisi della stabilità}.


\paragraph{Il problema della stabilità}

Consideriamo uno stato stazionario, o \emph{steady state},
$\mathbf{x}^*$, definito dalla condizione
\[
\mathbf{x}^*=F(\mathbf{x}^*).
\]
Se il sistema raggiunge tale stato, in assenza di perturbazioni vi
rimane:
\[
\mathbf{x}_t=\mathbf{x}^*
\quad\Rightarrow\quad
\mathbf{x}_{t+1}=\mathbf{x}^*.
\]

L'esistenza di uno steady state non implica, tuttavia, che il sistema
tenda verso di esso. La domanda rilevante è che cosa accade quando il
sistema parte da uno stato diverso, ma eventualmente vicino, a
$\mathbf{x}^*$.

In termini intuitivi, uno steady state è \emph{stabile} se piccole
perturbazioni dello stato iniziale non producono traiettorie che si
allontanano arbitrariamente dall'equilibrio.

Una proprietà più forte è la \emph{stabilità asintotica}. Uno steady
state è localmente asintoticamente stabile se, partendo
sufficientemente vicino a $\mathbf{x}^*$, la traiettoria non soltanto
rimane nelle sue vicinanze, ma converge nuovamente verso di esso:
\[
\lim_{t\rightarrow\infty}\mathbf{x}_t=\mathbf{x}^*.
\]

È utile distinguere quindi due idee:
\[
\text{stabilità}
\quad\Rightarrow\quad
\text{le perturbazioni non allontanano il sistema dall'equilibrio},
\]
mentre
\[
\text{stabilità asintotica}
\quad\Rightarrow\quad
\text{le perturbazioni vengono progressivamente riassorbite}.
\]

L'analisi della stabilità permette quindi di andare oltre la
simulazione di singole traiettorie. Invece di provare numericamente un
numero molto elevato di condizioni iniziali, cerchiamo di individuare
proprietà della legge di movimento che consentano di caratterizzare
il comportamento delle traiettorie.

Nel caso lineare scalare
\[
x_{t+1}=a x_t+b,
\]
tale comportamento dipende dal coefficiente $a$. Nel caso di un
sistema lineare
\[
\mathbf{x}_{t+1}=A\mathbf{x}_t+\mathbf{b},
\]
il ruolo svolto dal singolo coefficiente $a$ viene assunto dalla
matrice $A$.

L'analisi della stabilità dei sistemi multidimensionali conduce così
naturalmente allo studio degli autovalori della matrice di transizione.
In particolare, la posizione degli autovalori rispetto al cerchio
unitario permette di stabilire se le perturbazioni tendono a
scomparire, a persistere oppure ad amplificarsi nel tempo.

Nei modelli forward-looking, come vedremo, gli autovalori assumono
anche un ruolo ulteriore. La presenza di radici esterne al cerchio
unitario non implica necessariamente l'instabilità della soluzione:
le variabili forward-looking possono aggiustarsi in modo da eliminare
le componenti esplosive. Il numero delle radici stabili e instabili
diventa quindi rilevante anche per determinare l'esistenza e
l'unicità della soluzione.

Il calcolo degli autovalori diventa pertanto un problema
computazionale centrale nello studio dei sistemi dinamici discreti.
Per sistemi di piccole dimensioni gli autovalori possono essere
analizzati direttamente; per sistemi più grandi è invece necessario
ricorrere ad algoritmi numerici specifici, che verranno discussi in
seguito.

Per i sistemi ricorsivi la sequenza concettuale che seguiremo sarà
quindi:
\begin{itemize}
    \item condizione iniziale;
    \item traiettoria;
    \item variazione di condizioni iniziali e parametri;
    \item steady state;
    \item stabilità;
    \item autovalori;
    \item metodi numerici.
\end{itemize}

Per i sistemi forward-looking questa sequenza dovrà essere preceduta
dalla soluzione del problema intertemporale, utilizzando congiuntamente
le condizioni iniziali e le condizioni imposte sul comportamento futuro
della traiettoria.


\chapter{Modelli deterministici}

\section{Equazioni alle differenze lineari del primo ordine}

Il punto di partenza più semplice per lo studio dei modelli dinamici
discreti è costituito da una singola equazione alle differenze lineare
del primo ordine. La sua forma generale è
\[
x_{t+1}=a x_t+b,
\]
dove $x_t$ rappresenta lo stato del sistema al tempo $t$, mentre $a$ e
$b$ sono parametri.

L'equazione è detta \emph{del primo ordine} perché il valore della
variabile al periodo successivo dipende dal valore assunto nel periodo
corrente. È detta \emph{lineare} perché la relazione tra $x_{t+1}$ e
$x_t$ è lineare.

Per ottenere una particolare traiettoria è necessario specificare una
condizione iniziale:
\[
x_0=\bar{x}.
\]
Il problema
\[
\begin{cases}
x_{t+1}=a x_t+b,\\
x_0=\bar{x}
\end{cases}
\]
costituisce quindi un \emph{Initial Value Problem}.

Dal punto di vista computazionale, questo è uno dei problemi dinamici
più semplici. Una volta conosciuto $x_t$, il valore successivo può
essere calcolato direttamente:
\[
x_{t+1}=a x_t+b.
\]
Ripetendo l'operazione si ottiene
\[
x_0\longrightarrow x_1\longrightarrow x_2
\longrightarrow\cdots\longrightarrow x_T.
\]

Questa procedura prende il nome di \emph{forward iteration}. Non è
necessario risolvere alcun sistema simultaneo: a ogni periodo tutte le
informazioni necessarie per determinare lo stato successivo sono già
disponibili.


\subsection{Dalla condizione iniziale alla traiettoria}

Consideriamo, ad esempio, il modello
\[
x_{t+1}=0.8x_t+2
\]
e supponiamo che
\[
x_0=5.
\]
I primi valori della traiettoria sono
\[
x_1=0.8(5)+2=6,
\]
\[
x_2=0.8(6)+2=6.8,
\]
\[
x_3=0.8(6.8)+2=7.44,
\]
e così via.

Il computer esegue esattamente questa operazione, ma può ripeterla per
un numero molto elevato di periodi. Un semplice algoritmo può essere
rappresentato schematicamente come
\[
\begin{aligned}
&x[0]\leftarrow x_0,\\
&\text{per }t=0,\ldots,T-1:\\
&\qquad x[t+1]\leftarrow a\,x[t]+b.
\end{aligned}
\]

La semplicità dell'algoritmo non deve far pensare che il problema
dinamico sia privo di interesse. La simulazione permette infatti di
porre immediatamente alcune domande.

Che cosa accade se cambiamo $x_0$? Le diverse traiettorie tendono verso
lo stesso valore? Che cosa accade se modifichiamo $a$? La traiettoria
continua a convergere oppure può iniziare a oscillare o a divergere?

È utile distinguere chiaramente questi due esperimenti.

Se cambiamo la condizione iniziale mantenendo invariati $a$ e $b$,
consideriamo diversi IVP associati alla stessa legge di movimento. Se
cambiamo $a$ oppure $b$, modifichiamo invece la legge di movimento e
quindi il sistema dinamico stesso.


\subsection{Lo steady state}

Per studiare sistematicamente queste traiettorie è utile individuare
innanzitutto gli stati che rimangono invariati nel tempo.

Uno \emph{steady state}, o punto fisso, $x^*$ soddisfa
\[
x^*=a x^*+b.
\]
Se $a\neq 1$, possiamo quindi ottenere
\[
x^*=\frac{b}{1-a}.
\]

Nel precedente esempio,
\[
x_{t+1}=0.8x_t+2,
\]
lo steady state è
\[
x^*=\frac{2}{1-0.8}=10.
\]

Se $x_t=10$, infatti,
\[
x_{t+1}=0.8(10)+2=10.
\]

Lo steady state rappresenta quindi un valore che, una volta raggiunto,
viene riprodotto dalla legge di movimento.

Questo risultato, tuttavia, non ci dice ancora se una traiettoria che
parte da un valore diverso da $x^*$ tenderà verso lo steady state.


\subsection{La dinamica rispetto allo steady state}

Per studiare il comportamento della traiettoria è conveniente
considerare la deviazione dallo steady state:
\[
\widetilde{x}_t=x_t-x^*.
\]

Poiché
\[
x_{t+1}=a x_t+b
\]
e
\[
x^*=a x^*+b,
\]
sottraendo la seconda relazione dalla prima otteniamo
\[
x_{t+1}-x^*=a(x_t-x^*).
\]
Pertanto
\[
\widetilde{x}_{t+1}=a\widetilde{x}_t.
\]

Iterando questa relazione:
\[
\widetilde{x}_1=a\widetilde{x}_0,
\]
\[
\widetilde{x}_2=a^2\widetilde{x}_0,
\]
e, in generale,
\[
\widetilde{x}_t=a^t\widetilde{x}_0.
\]

Possiamo quindi scrivere la soluzione come
\[
x_t=x^*+a^t(x_0-x^*).
\]

Questa espressione è particolarmente importante. Essa separa la
traiettoria in due componenti:
\[
\underbrace{x^*}_{\text{steady state}}
+
\underbrace{a^t(x_0-x^*)}_{\text{deviazione dallo steady state}}.
\]

Il comportamento dinamico dipende quindi dal comportamento di $a^t$
quando $t$ cresce.


\subsection{Il ruolo del parametro dinamico}

Consideriamo inizialmente
\[
0<a<1.
\]
In questo caso
\[
a^t\rightarrow0
\]
e quindi
\[
x_t\rightarrow x^*.
\]
Inoltre, il segno della deviazione iniziale non cambia: la traiettoria
si avvicina progressivamente allo steady state senza alternare valori
al di sopra e al di sotto di esso.

Abbiamo quindi una \emph{convergenza monotona}.

Se invece
\[
-1<a<0,
\]
abbiamo ancora
\[
a^t\rightarrow0,
\]
ma il segno di $a^t$ cambia a ogni periodo. La traiettoria passa
alternativamente da un lato all'altro dello steady state, mentre
l'ampiezza delle oscillazioni si riduce progressivamente.

Abbiamo quindi una \emph{convergenza oscillatoria}.

Se
\[
a>1,
\]
il valore assoluto di $a^t$ cresce. Una deviazione iniziale dallo
steady state viene quindi progressivamente amplificata. La traiettoria
si allontana dallo steady state senza alternarne il lato.

Se
\[
a<-1,
\]
anche il valore assoluto della deviazione aumenta, ma il suo segno
cambia a ogni periodo. Otteniamo pertanto oscillazioni di ampiezza
crescente.

I casi limite $a=1$ e $a=-1$ richiedono un'attenzione particolare.
Nel primo caso le deviazioni non vengono ridotte:
\[
\widetilde{x}_{t+1}=\widetilde{x}_t.
\]
Nel secondo:
\[
\widetilde{x}_{t+1}=-\widetilde{x}_t,
\]
e il sistema continua a oscillare con ampiezza costante.

La condizione fondamentale per la convergenza verso lo steady state è
quindi
\[
|a|<1.
\]

In questo caso lo steady state è asintoticamente stabile.


\subsection{Stabilità e dipendenza dalle condizioni iniziali}

Il risultato precedente permette di comprendere meglio la differenza
tra simulazione e analisi della stabilità.

Consideriamo due condizioni iniziali $x_0^{(1)}$ e $x_0^{(2)}$. Le
corrispondenti traiettorie soddisfano
\[
x_t^{(1)}
=
x^*+a^t(x_0^{(1)}-x^*)
\]
e
\[
x_t^{(2)}
=
x^*+a^t(x_0^{(2)}-x^*).
\]

La loro distanza è
\[
x_t^{(1)}-x_t^{(2)}
=
a^t\left(x_0^{(1)}-x_0^{(2)}\right).
\]

Se
\[
|a|<1,
\]
allora
\[
\left|x_t^{(1)}-x_t^{(2)}\right|
\rightarrow0.
\]

Condizioni iniziali differenti producono quindi traiettorie differenti
nel breve periodo, ma l'influenza della condizione iniziale scompare
progressivamente.

Se invece
\[
|a|>1,
\]
piccole differenze iniziali vengono amplificate.

La stabilità consente pertanto di formulare una conclusione sul
comportamento di un'intera famiglia di traiettorie senza dover
simulare separatamente ogni possibile condizione iniziale.


\subsection{Un esempio economico: il modello cobweb}

Un'applicazione classica delle equazioni alle differenze è il modello
\emph{cobweb}, utilizzato per descrivere l'aggiustamento dinamico dei
prezzi in mercati nei quali l'offerta reagisce con ritardo alle
condizioni di mercato.

Consideriamo una funzione di domanda
\[
Q_t^D=\alpha-\beta P_t,
\qquad \beta>0,
\]
e supponiamo che la quantità offerta nel periodo $t$ sia stata decisa
sulla base del prezzo osservato nel periodo precedente:
\[
Q_t^S=\gamma+\delta P_{t-1},
\qquad \delta>0.
\]

L'ipotesi economica fondamentale è quindi che i produttori non possano
modificare immediatamente la quantità offerta. La produzione del
periodo $t$ è stata programmata utilizzando l'informazione disponibile
nel periodo $t-1$.

L'equilibrio di mercato richiede
\[
Q_t^D=Q_t^S.
\]
Pertanto
\[
\alpha-\beta P_t
=
\gamma+\delta P_{t-1}.
\]

Risolvendo rispetto a $P_t$ otteniamo
\[
P_t
=
\frac{\alpha-\gamma}{\beta}
-
\frac{\delta}{\beta}P_{t-1}.
\]

Il modello cobweb assume quindi esattamente la forma
\[
x_{t+1}=ax_t+b,
\]
con
\[
a=-\frac{\delta}{\beta}.
\]

Il segno negativo del coefficiente dinamico ha una conseguenza
immediata: quando il sistema converge, la convergenza avviene attraverso
oscillazioni.

La condizione di stabilità è
\[
\left|-\frac{\delta}{\beta}\right|<1,
\]
ossia
\[
\frac{\delta}{\beta}<1.
\]

Pertanto
\[
\delta<\beta.
\]

L'interpretazione economica è particolarmente intuitiva. Se l'offerta
reagisce al prezzo precedente meno intensamente di quanto la domanda
reagisca al prezzo corrente, le deviazioni dall'equilibrio tendono a
ridursi. Se invece la risposta dell'offerta è troppo forte, il
meccanismo di aggiustamento può amplificare le oscillazioni.

Il modello mostra quindi come una condizione matematica di stabilità
possa essere tradotta in una condizione sui parametri economici.


\subsection{Un esempio finanziario: aggiustamento verso il valore fondamentale}

La stessa struttura matematica può essere utilizzata per rappresentare
un semplice processo di aggiustamento di un prezzo finanziario verso
un valore fondamentale.

Indichiamo con $P_t$ il prezzo dell'attività finanziaria e con $P^*$
il suo valore fondamentale, che supponiamo per il momento costante.

Consideriamo la regola di aggiustamento
\[
P_{t+1}-P_t
=
\alpha(P^*-P_t),
\]
dove $\alpha$ misura la velocità con cui il prezzo reagisce alla
distanza dal valore fondamentale.

Se
\[
P_t<P^*,
\]
la differenza $P^*-P_t$ è positiva e il meccanismo tende ad aumentare
il prezzo. Se
\[
P_t>P^*,
\]
la differenza è negativa e il meccanismo tende invece a ridurlo.

Riscrivendo l'equazione otteniamo
\[
P_{t+1}
=
(1-\alpha)P_t+\alpha P^*.
\]

Ancora una volta abbiamo la forma generale
\[
x_{t+1}=ax_t+b,
\]
con
\[
a=1-\alpha
\]
e
\[
b=\alpha P^*.
\]

Lo steady state soddisfa
\[
P^*
=
(1-\alpha)P^*+\alpha P^*,
\]
ed è quindi precisamente il valore fondamentale.

Definiamo ora il mispricing
\[
m_t=P_t-P^*.
\]
La sua dinamica è
\[
m_{t+1}
=
(1-\alpha)m_t.
\]

La condizione di convergenza diventa pertanto
\[
|1-\alpha|<1,
\]
ossia
\[
0<\alpha<2.
\]

Il valore del parametro $\alpha$ determina inoltre il tipo di
aggiustamento.

Se
\[
0<\alpha<1,
\]
il prezzo converge monotonicamente verso il valore fondamentale.

Se
\[
1<\alpha<2,
\]
il prezzo reagisce più che proporzionalmente al mispricing corrente.
Esso supera quindi il valore fondamentale e la traiettoria converge
attraverso oscillazioni.

Se
\[
\alpha>2,
\]
l'overshooting è sufficientemente forte da produrre oscillazioni di
ampiezza crescente e il valore fondamentale diventa instabile.

Questo semplice modello non intende rappresentare in modo completo il
processo di formazione dei prezzi nei mercati finanziari. Il suo
obiettivo è mostrare come un meccanismo economico di aggiustamento possa
essere tradotto in un'equazione alle differenze e analizzato con gli
stessi strumenti utilizzati per il modello cobweb.


\subsection{Due modelli, lo stesso problema computazionale}

I due esempi considerati descrivono fenomeni molto diversi.

Nel modello cobweb la dinamica nasce dal ritardo con cui i produttori
adeguano l'offerta alle condizioni di mercato. Nel modello finanziario
la dinamica deriva invece da un meccanismo di correzione della distanza
tra prezzo osservato e valore fondamentale.

Dal punto di vista computazionale, tuttavia, entrambi conducono alla
stessa struttura:
\[
x_{t+1}=ax_t+b.
\]

In entrambi i casi dobbiamo:

\begin{enumerate}
    \item specificare i parametri;
    \item individuare la condizione iniziale;
    \item iterare la legge di movimento;
    \item costruire la traiettoria;
    \item determinare lo steady state;
    \item verificare se le deviazioni dallo steady state tendono a
    ridursi o ad amplificarsi.
\end{enumerate}

Questo esempio illustra una caratteristica fondamentale
dell'economia e della finanza computazionale. Modelli con
interpretazioni economiche molto diverse possono generare lo stesso
problema matematico e richiedere, di conseguenza, le stesse tecniche
numeriche.

Nel caso di una singola equazione lineare la dinamica è governata da
un solo coefficiente, $a$, e la condizione
\[
|a|<1
\]
permette di caratterizzare completamente la stabilità dello steady
state.

Nei modelli economici e finanziari più realistici, tuttavia, lo stato
del sistema è generalmente descritto da più variabili che interagiscono
tra loro. La legge di movimento assume allora la forma
\[
\mathbf{x}_{t+1}=A\mathbf{x}_t+\mathbf{b}.
\]

Il problema che affronteremo nella prossima sezione sarà quindi quello
di comprendere quale sia, per una matrice $A$, l'analogo della semplice
condizione $|a|<1$. Questo passaggio condurrà naturalmente allo studio
degli autovalori e alle tecniche computazionali utilizzate per
determinarli.

\section{Sistemi lineari di equazioni alle differenze}

Nella sezione precedente abbiamo considerato modelli dinamici descritti da
una sola variabile di stato:
\[
x_{t+1}=ax_t+b.
\]
In questo caso il comportamento dinamico dipende da un unico coefficiente,
$a$, e la condizione
\[
|a|<1
\]
garantisce la convergenza verso lo steady state.

Molti modelli economici e finanziari richiedono tuttavia più variabili di
stato. Produzione, prezzi, reddito, capitale, tassi di interesse e prezzi
delle attività possono evolvere congiuntamente e influenzarsi
reciprocamente.

Consideriamo, ad esempio, due variabili:
\[
x_{1,t}
\qquad\text{e}\qquad
x_{2,t}.
\]
Un sistema lineare del primo ordine può essere scritto come
\[
\begin{aligned}
x_{1,t+1}
&=
a_{11}x_{1,t}+a_{12}x_{2,t}+b_1,\\
x_{2,t+1}
&=
a_{21}x_{1,t}+a_{22}x_{2,t}+b_2.
\end{aligned}
\]

Utilizzando la notazione matriciale, definiamo
\[
\mathbf{x}_t=
\begin{bmatrix}
x_{1,t}\\
x_{2,t}
\end{bmatrix},
\qquad
A=
\begin{bmatrix}
a_{11}&a_{12}\\
a_{21}&a_{22}
\end{bmatrix},
\qquad
\mathbf{b}=
\begin{bmatrix}
b_1\\
b_2
\end{bmatrix}.
\]
Il sistema può allora essere rappresentato mediante una sola equazione:
\[
\mathbf{x}_{t+1}
=
A\mathbf{x}_t+\mathbf{b}.
\]

La matrice $A$ prende il nome di \emph{matrice di transizione}. I suoi
elementi descrivono il modo in cui lo stato corrente influenza quello del
periodo successivo.

Gli elementi sulla diagonale,
\[
a_{11},\qquad a_{22},
\]
misurano la dipendenza di ciascuna variabile dal proprio valore precedente,
mentre gli elementi fuori diagonale,
\[
a_{12},\qquad a_{21},
\]
rappresentano le interazioni dinamiche tra le variabili.


\subsection{L'Initial Value Problem}

Anche nel caso multidimensionale è necessario specificare una condizione
iniziale:
\[
\mathbf{x}_0=\bar{\mathbf{x}}.
\]
L'Initial Value Problem è quindi
\[
\begin{cases}
\mathbf{x}_{t+1}=A\mathbf{x}_t+\mathbf{b},\\
\mathbf{x}_0=\bar{\mathbf{x}}.
\end{cases}
\]

La forward iteration procede esattamente come nel caso scalare:
\[
\mathbf{x}_1=A\mathbf{x}_0+\mathbf{b},
\]
\[
\mathbf{x}_2=A\mathbf{x}_1+\mathbf{b},
\]
e così via.

L'algoritmo può essere rappresentato schematicamente come
\[
\begin{aligned}
&\mathbf{x}[0]\leftarrow\mathbf{x}_0,\\
&\text{per }t=0,\ldots,T-1:\\
&\qquad
\mathbf{x}[t+1]\leftarrow
A\mathbf{x}[t]+\mathbf{b}.
\end{aligned}
\]

Dal punto di vista concettuale, quindi, il passaggio da una variabile a
molte variabili non modifica la logica della simulazione. Cambia invece
la complessità della dinamica che può essere generata.


\subsection{Lo steady state}

Uno steady state $\mathbf{x}^*$ soddisfa
\[
\mathbf{x}^*
=
A\mathbf{x}^*+\mathbf{b}.
\]
Riordinando:
\[
(I-A)\mathbf{x}^*
=
\mathbf{b}.
\]

Se la matrice $I-A$ è non singolare, lo steady state è unico e può essere
ottenuto risolvendo il sistema lineare
\[
(I-A)\mathbf{x}^*=\mathbf{b}.
\]

Formalmente possiamo scrivere
\[
\mathbf{x}^*
=
(I-A)^{-1}\mathbf{b},
\]
ma, come già osservato nello studio dei sistemi lineari, dal punto di vista
computazionale è generalmente preferibile risolvere direttamente il sistema
piuttosto che calcolare esplicitamente l'inversa.

Ritroviamo quindi una tecnica computazionale già studiata:
\[
\boxed{
\text{calcolo dello steady state}
\longrightarrow
\text{soluzione di un sistema lineare}
}.
\]


\subsection{La dinamica rispetto allo steady state}

Definiamo la deviazione dallo steady state:
\[
\widetilde{\mathbf{x}}_t
=
\mathbf{x}_t-\mathbf{x}^*.
\]

Poiché
\[
\mathbf{x}_{t+1}=A\mathbf{x}_t+\mathbf{b}
\]
e
\[
\mathbf{x}^*=A\mathbf{x}^*+\mathbf{b},
\]
sottraendo le due equazioni otteniamo
\[
\widetilde{\mathbf{x}}_{t+1}
=
A\widetilde{\mathbf{x}}_t.
\]

Iterando:
\[
\widetilde{\mathbf{x}}_1
=
A\widetilde{\mathbf{x}}_0,
\]
\[
\widetilde{\mathbf{x}}_2
=
A^2\widetilde{\mathbf{x}}_0,
\]
e, in generale,
\[
\widetilde{\mathbf{x}}_t
=
A^t\widetilde{\mathbf{x}}_0.
\]

La soluzione può pertanto essere scritta come
\[
\mathbf{x}_t
=
\mathbf{x}^*
+
A^t(\mathbf{x}_0-\mathbf{x}^*).
\]

Questa espressione costituisce l'analogo multidimensionale della soluzione
ottenuta nel caso scalare:
\[
x_t=x^*+a^t(x_0-x^*).
\]

Il problema della stabilità può quindi essere formulato in modo molto
semplice. Nel caso scalare ci siamo chiesti se
\[
a^t\rightarrow0.
\]
Nel caso multidimensionale dobbiamo invece chiederci se
\[
A^t\rightarrow0.
\]


\subsection{Perché non basta guardare gli elementi di $A$?}

Nel caso scalare esiste un unico coefficiente dinamico. In un sistema
multidimensionale, invece, le variabili interagiscono tra loro.

Consideriamo, ad esempio,
\[
A=
\begin{bmatrix}
a_{11}&a_{12}\\
a_{21}&a_{22}
\end{bmatrix}.
\]

Non è sufficiente analizzare separatamente $a_{11}$ e $a_{22}$, perché
l'evoluzione di $x_{1,t}$ dipende anche da $x_{2,t}$ e viceversa.

Una variazione di $x_{1,t}$ può modificare $x_{2,t+1}$ attraverso
$a_{21}$; la variazione di $x_{2,t+1}$ può successivamente tornare a
influenzare $x_{1,t+2}$ attraverso $a_{12}$.

Si generano quindi meccanismi di retroazione dinamica che dipendono
dall'intera matrice $A$.

Abbiamo bisogno di uno strumento che permetta di sintetizzare queste
interazioni.


\subsection{Autovalori e autovettori}

Consideriamo un vettore non nullo $\mathbf{v}$ tale che
\[
A\mathbf{v}=\lambda\mathbf{v}.
\]
Il numero $\lambda$ è un \emph{autovalore} di $A$ e $\mathbf{v}$ è il
corrispondente \emph{autovettore}.

L'importanza dinamica di questa relazione è immediata. Se una perturbazione
si trova nella direzione dell'autovettore $\mathbf{v}$, dopo un periodo
avremo
\[
A\mathbf{v}=\lambda\mathbf{v}.
\]
Dopo due periodi:
\[
A^2\mathbf{v}
=
A(\lambda\mathbf{v})
=
\lambda^2\mathbf{v}.
\]
Dopo $t$ periodi:
\[
A^t\mathbf{v}
=
\lambda^t\mathbf{v}.
\]

L'autovalore svolge quindi, lungo la direzione individuata
dall'autovettore, esattamente il ruolo che il coefficiente $a$ svolgeva nel
modello scalare.

Se
\[
|\lambda|<1,
\]
la componente associata a quell'autovalore tende ad attenuarsi.

Se
\[
|\lambda|>1,
\]
essa tende invece ad amplificarsi.

Se l'autovalore è negativo o complesso, la dinamica può inoltre presentare
oscillazioni.


\subsection{La condizione di stabilità}

Un sistema lineare discreto è asintoticamente stabile se tutte le componenti
della deviazione dallo steady state tendono a zero.

Questo richiede che tutti gli autovalori della matrice $A$ abbiano modulo
inferiore a uno:
\[
|\lambda_i|<1
\qquad
\text{per ogni }i.
\]

Definiamo il \emph{raggio spettrale} della matrice come
\[
\rho(A)
=
\max_i|\lambda_i|.
\]

La condizione di stabilità può quindi essere espressa in forma compatta:
\[
\boxed{\rho(A)<1}.
\]

Il risultato generalizza direttamente quello ottenuto nel caso scalare:
\[
|a|<1
\quad\longrightarrow\quad
\rho(A)<1.
\]

Il cerchio di raggio unitario nel piano complesso prende il nome di
\emph{cerchio unitario}. Per un sistema dinamico discreto lineare, la
stabilità asintotica richiede quindi che tutti gli autovalori si trovino
all'interno del cerchio unitario.


\subsection{Un esempio economico: un duopolio dinamico}

Consideriamo ora un mercato nel quale operano due imprese. Ciascuna impresa
sceglie la quantità da produrre nel periodo successivo reagendo alla
quantità osservata dell'altra impresa nel periodo corrente.

Una semplice rappresentazione è
\[
q_{1,t+1}
=
\alpha_1-\beta_1 q_{2,t},
\]
\[
q_{2,t+1}
=
\alpha_2-\beta_2 q_{1,t},
\]
con
\[
\beta_1>0,
\qquad
\beta_2>0.
\]

L'interpretazione è che una maggiore produzione del concorrente induce
ciascuna impresa a ridurre la propria produzione nel periodo successivo.

Definiamo
\[
\mathbf{q}_t=
\begin{bmatrix}
q_{1,t}\\
q_{2,t}
\end{bmatrix}.
\]
Il sistema diventa
\[
\mathbf{q}_{t+1}
=
\begin{bmatrix}
0&-\beta_1\\
-\beta_2&0
\end{bmatrix}
\mathbf{q}_t
+
\begin{bmatrix}
\alpha_1\\
\alpha_2
\end{bmatrix}.
\]

La matrice di transizione è quindi
\[
A=
\begin{bmatrix}
0&-\beta_1\\
-\beta_2&0
\end{bmatrix}.
\]

Gli autovalori soddisfano
\[
\det(A-\lambda I)=0.
\]
Pertanto
\[
\det
\begin{bmatrix}
-\lambda&-\beta_1\\
-\beta_2&-\lambda
\end{bmatrix}
=0,
\]
da cui
\[
\lambda^2-\beta_1\beta_2=0.
\]

Gli autovalori sono quindi
\[
\lambda_1=\sqrt{\beta_1\beta_2},
\qquad
\lambda_2=-\sqrt{\beta_1\beta_2}.
\]

Il raggio spettrale è
\[
\rho(A)=\sqrt{\beta_1\beta_2}.
\]

La condizione di stabilità è pertanto
\[
\sqrt{\beta_1\beta_2}<1,
\]
equivalentemente
\[
\beta_1\beta_2<1.
\]

La stabilità dipende quindi dall'interazione tra le reazioni delle due
imprese. Anche se ciascuna equazione appare semplice, le decisioni delle
imprese generano un circuito di retroazione:
\[
q_{1,t}
\longrightarrow
q_{2,t+1}
\longrightarrow
q_{1,t+2}
\longrightarrow\cdots.
\]

Se il prodotto delle intensità di reazione è sufficientemente piccolo, le
perturbazioni vengono progressivamente attenuate. Se le reazioni sono
troppo forti, il meccanismo di risposta reciproca può invece amplificare le
deviazioni dall'equilibrio.


\subsection{Un esempio finanziario: prezzo e momentum}

Consideriamo ora un semplice modello nel quale la dinamica del prezzo di
un'attività finanziaria dipende sia dalla distanza dal valore fondamentale
sia dal movimento recente del prezzo.

Indichiamo con $P^*$ il valore fondamentale e definiamo il mispricing:
\[
m_t=P_t-P^*.
\]

Introduciamo inoltre una variabile $z_t$ che rappresenta una componente di
momentum, ossia una misura sintetica della recente tendenza del prezzo.

Consideriamo il sistema
\[
m_{t+1}
=
(1-\alpha)m_t+\gamma z_t,
\]
\[
z_{t+1}
=
\delta m_t+\phi z_t.
\]

Il parametro $\alpha$ rappresenta la forza del meccanismo di ritorno verso
il valore fondamentale, mentre $\gamma$, $\delta$ e $\phi$ descrivono la
trasmissione e la persistenza della componente di momentum.

Definendo
\[
\mathbf{x}_t=
\begin{bmatrix}
m_t\\
z_t
\end{bmatrix},
\]
otteniamo
\[
\mathbf{x}_{t+1}
=
A\mathbf{x}_t,
\]
con
\[
A=
\begin{bmatrix}
1-\alpha&\gamma\\
\delta&\phi
\end{bmatrix}.
\]

Lo steady state naturale è
\[
\mathbf{x}^*
=
\begin{bmatrix}
0\\
0
\end{bmatrix},
\]
che corrisponde a
\[
P_t=P^*
\]
e all'assenza di momentum.

La dinamica dipende tuttavia dall'interazione tra i due meccanismi.

Il ritorno verso il fondamentale tende a correggere il mispricing, mentre
il momentum può propagare nel tempo i movimenti precedenti del prezzo.

Per determinare quale dei due meccanismi prevalga non è sufficiente
considerare separatamente i parametri. Occorre analizzare gli autovalori
della matrice
\[
A=
\begin{bmatrix}
1-\alpha&\gamma\\
\delta&\phi
\end{bmatrix}.
\]

Se
\[
\rho(A)<1,
\]
una perturbazione del prezzo o del momentum viene progressivamente
riassorbita e il sistema ritorna verso
\[
(m,z)=(0,0).
\]

Se invece
\[
\rho(A)>1,
\]
esiste almeno una componente della dinamica che tende ad amplificarsi.

L'esempio mostra quindi come l'interazione tra meccanismi di correzione e
meccanismi di propagazione possa essere studiata attraverso le proprietà
spettrali della matrice di transizione.


\subsection{Dalla simulazione all'analisi spettrale}

Il duopolio e il modello finanziario descrivono fenomeni differenti, ma
conducono allo stesso problema computazionale:
\[
\mathbf{x}_{t+1}
=
A\mathbf{x}_t+\mathbf{b}.
\]

Una possibile strategia consiste nel simulare il modello per numerosi
periodi e osservare se le variabili sembrano convergere.

Questo approccio è utile per visualizzare la dinamica, ma presenta un
limite importante. Una simulazione riguarda una specifica condizione
iniziale e un particolare orizzonte temporale. Una traiettoria che sembra
stabile per un certo numero di periodi non costituisce, da sola, una prova
della stabilità del sistema.

L'analisi degli autovalori consente invece di caratterizzare direttamente
la dinamica asintotica.

Per un sistema lineare discreto:
\[
\boxed{
\rho(A)<1
\quad\Longleftrightarrow\quad
\text{stabilità asintotica}
}.
\]

Il problema della stabilità si trasforma quindi in un problema di
\emph{eigenvalue computation}.


\subsection{Il calcolo degli autovalori come problema computazionale}

Per matrici molto piccole possiamo determinare gli autovalori risolvendo
l'equazione caratteristica
\[
\det(A-\lambda I)=0.
\]

Questa procedura è utile dal punto di vista teorico e permette di
comprendere la definizione di autovalore. Non rappresenta però il metodo
generale utilizzato dal computer per matrici di dimensione elevata.

Già per una matrice $n\times n$, l'equazione caratteristica è un polinomio
di grado $n$. All'aumentare della dimensione del sistema, costruire
esplicitamente tale polinomio e determinarne successivamente le radici
diventa una strategia numericamente poco conveniente.

Gli algoritmi numerici operano invece direttamente sulla matrice.

Per matrici dense di dimensione moderata, uno degli strumenti fondamentali
per il calcolo dello spettro è l'\emph{algoritmo QR}. L'idea generale
consiste nel costruire una successione di matrici simili ad $A$ che
conservano gli stessi autovalori ma assumono progressivamente una forma
dalla quale questi possono essere estratti più facilmente.

Quando i sistemi diventano molto grandi, tuttavia, può non essere
conveniente calcolare tutti gli autovalori.


\subsection{Sistemi grandi e matrici sparse}

Molti modelli economici e finanziari possono generare matrici di dimensione
elevata. Inoltre, ogni variabile interagisce spesso soltanto con una parte
limitata delle altre variabili.

La matrice di transizione può quindi essere \emph{sparsa}: molti dei suoi
elementi sono uguali a zero.

In questi casi emergono due considerazioni computazionali.

La prima riguarda la memoria. Memorizzare tutti gli $n^2$ elementi di una
matrice molto grande è inefficiente quando soltanto una piccola frazione di
essi è diversa da zero. Le rappresentazioni sparse memorizzano quindi
soltanto gli elementi non nulli e la loro posizione.

La seconda riguarda il calcolo degli autovalori. Per studiare la stabilità
non siamo necessariamente interessati all'intero spettro. La condizione
\[
\rho(A)<1
\]
dipende soprattutto dagli autovalori di modulo maggiore.

Diventa quindi naturale cercare direttamente gli autovalori dominanti,
evitando di calcolare quelli che non sono rilevanti per il problema.


\subsection{Metodi iterativi per gli autovalori dominanti}

Il più semplice esempio è il \emph{power method}. Partendo da un vettore
iniziale $\mathbf{v}_0$, si costruisce la successione
\[
\mathbf{v}_{k+1}
=
\frac{A\mathbf{v}_k}
{\|A\mathbf{v}_k\|}.
\]

Sotto opportune condizioni, la direzione di $\mathbf{v}_k$ tende verso
l'autovettore associato all'autovalore dominante.

Il metodo mostra un principio computazionale importante: se siamo
interessati soltanto alla parte dello spettro rilevante per la stabilità,
può non essere necessario risolvere l'intero problema agli autovalori.

Per problemi più complessi vengono utilizzati metodi basati sui
\emph{Krylov subspaces}, tra i quali il metodo di Arnoldi per matrici
generali.

L'idea è costruire spazi della forma
\[
\mathcal{K}_m(A,\mathbf{v})
=
\operatorname{span}
\left\{
\mathbf{v},
A\mathbf{v},
A^2\mathbf{v},
\ldots,
A^{m-1}\mathbf{v}
\right\},
\]
all'interno dei quali è possibile ottenere approssimazioni degli
autovalori rilevanti senza trattare esplicitamente l'intero problema
spettrale.

La scelta del metodo dipende quindi dalla struttura del problema:

\begin{itemize}
    \item per sistemi molto piccoli, il calcolo diretto permette anche
    un'analisi teorica;
    \item per matrici dense di dimensione moderata, si utilizzano
    algoritmi numerici per il problema completo agli autovalori;
    \item per matrici grandi e sparse, sono spesso preferibili metodi
    iterativi che determinano soltanto una parte dello spettro.
\end{itemize}


\subsection{Una nuova prospettiva sulla stabilità}

Il passaggio dalla singola equazione ai sistemi mostra bene il ruolo delle
tecniche computazionali nell'analisi dinamica.

Nel caso scalare avevamo
\[
x_{t+1}=ax_t+b
\]
e la stabilità dipendeva da
\[
|a|<1.
\]

Nel caso multidimensionale abbiamo
\[
\mathbf{x}_{t+1}
=
A\mathbf{x}_t+\mathbf{b}
\]
e la condizione diventa
\[
\rho(A)<1.
\]

La struttura logica rimane quindi la stessa:
\[
\boxed{
\text{legge di movimento}
\longrightarrow
\text{traiettorie}
\longrightarrow
\text{steady state}
\longrightarrow
\text{stabilità}
\longrightarrow
\text{autovalori}
}.
\]

Ciò che cambia è il problema computazionale. Al crescere del numero delle
variabili, la simulazione richiede ripetute operazioni matrice-vettore,
il calcolo dello steady state richiede la soluzione di un sistema lineare
e l'analisi della stabilità richiede informazioni sullo spettro della
matrice di transizione.

Finora abbiamo però considerato sistemi nei quali, dato lo stato corrente,
lo stato successivo è immediatamente determinabile:
\[
\mathbf{x}_{t+1}
=
A\mathbf{x}_t+\mathbf{b}.
\]

Non tutti i modelli economici e finanziari possiedono questa struttura.
Quando alcune variabili del nuovo periodo sono determinate
simultaneamente, il modello può presentarsi invece in forma implicita:
\[
G(\mathbf{x}_{t+1},\mathbf{x}_t)=0.
\]

Prima di poter effettuare la forward iteration sarà allora necessario
risolvere, a ogni periodo, un ulteriore problema numerico. Questo sarà
l'oggetto della prossima sezione.

\section{Sistemi dinamici in forma implicita}

Finora abbiamo considerato sistemi dinamici nei quali lo stato del periodo
successivo è espresso direttamente in funzione dello stato corrente:
\[
\mathbf{x}_{t+1}=F(\mathbf{x}_t).
\]

Questa rappresentazione è particolarmente conveniente perché consente di
procedere direttamente mediante forward iteration. Una volta noto
$\mathbf{x}_t$, il computer valuta la funzione $F$ e ottiene
$\mathbf{x}_{t+1}$.

Molti modelli economici e finanziari non vengono tuttavia formulati
direttamente in questa forma. In particolare, alcune variabili riferite allo
stesso periodo possono essere determinate simultaneamente.

Una rappresentazione più generale è
\[
G(\mathbf{x}_{t+1},\mathbf{x}_t)=0.
\]

In questo caso, dato $\mathbf{x}_t$, lo stato successivo non è ottenuto
mediante la semplice valutazione di una funzione. Occorre invece risolvere
il sistema
\[
G(\mathbf{x}_{t+1},\mathbf{x}_t)=0
\]
rispetto a $\mathbf{x}_{t+1}$.

Dal punto di vista computazionale, la dinamica assume quindi la struttura
\[
\mathbf{x}_t
\longrightarrow
\boxed{\text{soluzione del sistema}}
\longrightarrow
\mathbf{x}_{t+1}.
\]

La procedura deve essere ripetuta per ciascun periodo della simulazione.


\subsection{Forma esplicita e forma implicita}

È utile confrontare direttamente i due casi.

In un sistema esplicito:
\[
\mathbf{x}_{t+1}=F(\mathbf{x}_t),
\]
l'algoritmo può essere rappresentato schematicamente come
\[
\begin{aligned}
&\mathbf{x}[0]\leftarrow\mathbf{x}_0,\\
&\text{per }t=0,\ldots,T-1:\\
&\qquad
\mathbf{x}[t+1]\leftarrow
F(\mathbf{x}[t]).
\end{aligned}
\]

In un sistema implicito:
\[
G(\mathbf{x}_{t+1},\mathbf{x}_t)=0,
\]
l'algoritmo diventa invece
\[
\begin{aligned}
&\mathbf{x}[0]\leftarrow\mathbf{x}_0,\\
&\text{per }t=0,\ldots,T-1:\\
&\qquad
\text{risolvere }
G(\mathbf{x}[t+1],\mathbf{x}[t])=0
\\
&\qquad
\text{rispetto a }\mathbf{x}[t+1].
\end{aligned}
\]

La differenza è quindi importante: nel secondo caso una singola
transizione temporale incorpora a sua volta un problema numerico.


\subsection{Il caso lineare}

Consideriamo un sistema implicito lineare:
\[
A\mathbf{x}_{t+1}
=
B\mathbf{x}_t+\mathbf{c}.
\]

Se $A$ è non singolare, dal punto di vista matematico possiamo scrivere
\[
\mathbf{x}_{t+1}
=
A^{-1}B\mathbf{x}_t
+
A^{-1}\mathbf{c}.
\]

Il sistema può quindi essere ricondotto formalmente alla forma ricorsiva
\[
\mathbf{x}_{t+1}
=
C\mathbf{x}_t+\mathbf{d},
\]
dove
\[
C=A^{-1}B
\]
e
\[
\mathbf{d}=A^{-1}\mathbf{c}.
\]

Questa trasformazione è utile per analizzare teoricamente il sistema.

Dal punto di vista computazionale, tuttavia, non è generalmente conveniente
calcolare esplicitamente $A^{-1}$. A ogni periodo è preferibile risolvere
direttamente
\[
A\mathbf{x}_{t+1}
=
B\mathbf{x}_t+\mathbf{c}.
\]

Ritroviamo quindi una delle tecniche già studiate nella parte precedente
del corso:
\[
\boxed{
\text{sistema dinamico implicito lineare}
\longrightarrow
\text{linear solver a ogni periodo}
}.
\]


\subsection{Quando la forma ricorsiva esiste}

La presenza di simultaneità non significa necessariamente che il modello
debba essere risolto numericamente a ogni periodo.

In alcuni casi le equazioni possono essere manipolate algebricamente in
modo da ottenere direttamente
\[
\mathbf{x}_{t+1}=F(\mathbf{x}_t).
\]

Quando ciò è possibile, la simultaneità appartiene alla formulazione
originaria del modello, ma può essere eliminata prima della simulazione.

È utile quindi distinguere tra:

\begin{itemize}
    \item sistemi impliciti che possono essere trasformati analiticamente
    in forma ricorsiva;
    \item sistemi per i quali lo stato successivo deve essere determinato
    numericamente a ogni periodo.
\end{itemize}

Nel caso lineare, la possibilità di ottenere una forma ricorsiva dipende
in particolare dalla non singolarità della matrice che moltiplica le
variabili del nuovo periodo.


\subsection{Il caso non lineare}

Consideriamo ora
\[
G(\mathbf{x}_{t+1},\mathbf{x}_t)=0,
\]
dove $G$ è non lineare.

In generale potrebbe non essere possibile isolare analiticamente
$\mathbf{x}_{t+1}$.

Dato $\mathbf{x}_t$, dobbiamo quindi considerare
\[
H(\mathbf{y})
=
G(\mathbf{y},\mathbf{x}_t)
\]
e trovare un vettore $\mathbf{y}$ tale che
\[
H(\mathbf{y})=0.
\]

Una volta ottenuta la soluzione,
\[
\mathbf{x}_{t+1}=\mathbf{y},
\]
possiamo avanzare al periodo successivo.

La simulazione dinamica incorpora quindi ripetutamente un problema di
root finding:
\[
\boxed{
\text{dinamica}
+
\text{soluzione di sistemi non lineari}
}.
\]

Possono essere utilizzati, ad esempio, il metodo di Newton e le sue
varianti, già introdotti nella parte dedicata ai sistemi non lineari.


\subsection{Esistenza e unicità della transizione}

La possibilità di scrivere
\[
\mathbf{x}_{t+1}=F(\mathbf{x}_t)
\]
presuppone che, dato lo stato corrente, lo stato successivo sia ben
definito.

In un sistema implicito questo non è automatico.

L'equazione
\[
G(\mathbf{x}_{t+1},\mathbf{x}_t)=0
\]
può avere:

\begin{itemize}
    \item una soluzione;
    \item più soluzioni;
    \item nessuna soluzione.
\end{itemize}

Ad esempio,
\[
x_{t+1}^2=x_t
\]
per $x_t>0$ ammette due soluzioni:
\[
x_{t+1}=+\sqrt{x_t}
\]
e
\[
x_{t+1}=-\sqrt{x_t}.
\]

Non esiste quindi una singola legge di movimento senza introdurre
un'ulteriore regola di selezione.

Se invece
\[
x_t<0,
\]
non esiste alcuna soluzione reale.

Nei modelli dinamici impliciti diventano quindi rilevanti non soltanto il
calcolo numerico, ma anche i problemi di esistenza e unicità della
transizione.


\subsection{Una condizione locale: il teorema della funzione implicita}

Consideriamo
\[
G(\mathbf{x}_{t+1},\mathbf{x}_t)=0.
\]

Se, in un punto di interesse, la matrice Jacobiana rispetto alle variabili
future,
\[
\frac{\partial G}
{\partial\mathbf{x}_{t+1}},
\]
è non singolare, il teorema della funzione implicita garantisce, sotto
opportune condizioni di regolarità, che localmente sia possibile esprimere
le variabili future come funzione di quelle correnti:
\[
\mathbf{x}_{t+1}=F(\mathbf{x}_t).
\]

Il risultato generalizza al caso non lineare la condizione di invertibilità
incontrata per i sistemi lineari.

Dal punto di vista computazionale è importante perché suggerisce che,
almeno localmente, il problema implicito possiede una soluzione ben
definita che può essere ricercata numericamente.


\subsection{Un esempio economico: il modello SP-DG}

Un esempio di simultaneità intra-periodale è fornito dal modello SP-DG.

Consideriamo, in forma semplificata, una relazione di offerta aggregata:
\[
\pi_t
=
\pi_t^e+\alpha\widehat{Y}_t+z_t^s,
\]
dove $\pi_t$ rappresenta l'inflazione, $\pi_t^e$ l'inflazione attesa,
$\widehat{Y}_t$ l'output gap e $z_t^s$ uno shock di offerta.

La relazione di domanda aggregata può essere rappresentata come
\[
\widehat{Y}_t
=
\widehat{x}_t-\pi_t+\widehat{Y}_{t-1}-z_t^d,
\]
dove $\widehat{x}_t$ rappresenta una componente esogena della domanda e
$z_t^d$ uno shock di domanda.

Le due equazioni determinano simultaneamente
\[
\pi_t
\qquad\text{e}\qquad
\widehat{Y}_t.
\]

Non possiamo quindi calcolare prima $\pi_t$ senza conoscere
$\widehat{Y}_t$, ma non possiamo neppure calcolare
$\widehat{Y}_t$ senza conoscere $\pi_t$.

Il sistema è quindi, nella sua formulazione originaria, simultaneo.


\subsection{Eliminare la simultaneità nel modello SP-DG}

Sostituendo
\[
\pi_t
=
\pi_t^e+\alpha\widehat{Y}_t+z_t^s
\]
nella relazione di domanda otteniamo
\[
\widehat{Y}_t
=
\widehat{x}_t
-
\pi_t^e
-
\alpha\widehat{Y}_t
-
z_t^s
+
\widehat{Y}_{t-1}
-
z_t^d.
\]

Raccogliendo $\widehat{Y}_t$:
\[
(1+\alpha)\widehat{Y}_t
=
\widehat{x}_t
-
\pi_t^e
+
\widehat{Y}_{t-1}
-
z_t^s
-
z_t^d.
\]

Se
\[
1+\alpha\neq0,
\]
possiamo scrivere
\[
\widehat{Y}_t
=
\frac{
\widehat{x}_t
-
\pi_t^e
+
\widehat{Y}_{t-1}
-
z_t^s
-
z_t^d
}
{1+\alpha}.
\]

Una volta determinato $\widehat{Y}_t$, possiamo calcolare
\[
\pi_t
=
\pi_t^e+\alpha\widehat{Y}_t+z_t^s.
\]

La simultaneità è stata quindi eliminata analiticamente.

Il modello può essere simulato ricorsivamente dopo aver effettuato questa
trasformazione.


\subsection{Aspettative adattive nel modello SP-DG}

Nel modello considerato, le aspettative di inflazione sono adattive:
\[
\pi_t^e
=
\lambda\pi_{t-1}
+
(1-\lambda)\pi_{t-1}^e.
\]

L'inflazione attesa nel periodo $t$ dipende quindi esclusivamente da
quantità già note al termine del periodo precedente.

La sequenza computazionale può essere rappresentata come
\[
(\pi_{t-1},\pi_{t-1}^e,\widehat{Y}_{t-1})
\]
\[
\Downarrow
\]
\[
\pi_t^e
\]
\[
\Downarrow
\]
\[
\widehat{Y}_t
\]
\[
\Downarrow
\]
\[
\pi_t.
\]

Dopo aver risolto la simultaneità intra-periodale, il modello diventa
quindi pienamente ricorsivo.

Questo aspetto sarà importante quando confronteremo le aspettative
adattive con formulazioni forward-looking: nel primo caso il modello può
essere simulato procedendo periodo per periodo; nel secondo il problema
può richiedere una soluzione intertemporale preliminare.


\subsection{Un esempio finanziario: prezzo e rendimento di un'obbligazione}

Consideriamo ora un semplice esempio finanziario nel quale il rendimento
di un'obbligazione e il suo prezzo devono essere determinati congiuntamente.

Supponiamo, per semplicità, che un'obbligazione zero coupon paghi un valore
nominale $F$ dopo $n$ periodi.

Il prezzo e il rendimento a scadenza sono legati dalla relazione
\[
P_{t+1}
=
\frac{F}
{(1+y_{t+1})^n}.
\]

Supponiamo inoltre che il rendimento del periodo successivo sia determinato
da una relazione di mercato della forma
\[
y_{t+1}
=
\bar{y}
+
\gamma(P_{t+1}-P_t),
\]
dove $\bar{y}$ rappresenta un livello di riferimento e $\gamma$ misura la
risposta del rendimento alla variazione del prezzo.

Le due variabili
\[
P_{t+1}
\qquad\text{e}\qquad
y_{t+1}
\]
devono essere determinate simultaneamente.

Il sistema è
\[
\begin{cases}
P_{t+1}
=
\dfrac{F}{(1+y_{t+1})^n},\\[3mm]
y_{t+1}
=
\bar{y}
+
\gamma(P_{t+1}-P_t).
\end{cases}
\]

Dato $P_t$, non possiamo calcolare direttamente $P_{t+1}$ perché abbiamo
bisogno di $y_{t+1}$; allo stesso tempo, per determinare $y_{t+1}$ abbiamo
bisogno di $P_{t+1}$.

Si tratta quindi di un sistema dinamico implicito.


\subsection{Riduzione a un problema di root finding}

Sostituendo la seconda equazione nella prima otteniamo
\[
P_{t+1}
=
\frac{F}
{
\left[
1+\bar{y}
+
\gamma(P_{t+1}-P_t)
\right]^n
}.
\]

Portando tutti i termini a sinistra:
\[
H(P_{t+1};P_t)
=
P_{t+1}
-
\frac{F}
{
\left[
1+\bar{y}
+
\gamma(P_{t+1}-P_t)
\right]^n
}
=
0.
\]

Dato $P_t$, il valore di $P_{t+1}$ deve quindi essere determinato
numericamente risolvendo
\[
H(P_{t+1};P_t)=0.
\]

Una volta ottenuto $P_{t+1}$, possiamo calcolare
\[
y_{t+1}
=
\bar{y}
+
\gamma(P_{t+1}-P_t).
\]

L'algoritmo assume quindi la forma
\[
P_t
\longrightarrow
\boxed{\text{root finding}}
\longrightarrow
P_{t+1}
\longrightarrow
y_{t+1}.
\]

Questo esempio mostra come una relazione finanziaria non lineare possa
rendere implicita una transizione dinamica e richiedere la soluzione
numerica di un'equazione a ogni periodo.


\subsection{Il ruolo della soluzione iniziale nei metodi iterativi}

Quando a ogni periodo viene utilizzato un metodo iterativo, come Newton,
occorre specificare anche un valore iniziale per il solver.

Nel problema
\[
H(P_{t+1};P_t)=0,
\]
una scelta naturale può essere utilizzare
\[
P_t
\]
come valore iniziale per la ricerca di $P_{t+1}$.

Questa strategia sfrutta una caratteristica comune dei modelli dinamici:
quando la variazione tra periodi consecutivi non è eccessivamente grande,
lo stato corrente può costituire una buona approssimazione iniziale dello
stato successivo.

Il problema dinamico genera quindi anche una sequenza naturale di
\emph{warm starts} per il metodo numerico.


\subsection{Errore del solver e dinamica}

Nei sistemi espliciti:
\[
\mathbf{x}_{t+1}=F(\mathbf{x}_t),
\]
una volta specificata la legge di movimento, la transizione viene
calcolata direttamente, salvo gli ordinari errori numerici di
rappresentazione.

Nei sistemi impliciti risolti numericamente compare invece una nuova fonte
di approssimazione.

Supponiamo che il solver produca
\[
\widehat{\mathbf{x}}_{t+1}
\]
tale che
\[
G(\widehat{\mathbf{x}}_{t+1},\mathbf{x}_t)
\approx0.
\]

Il valore approssimato viene poi utilizzato come stato per il periodo
successivo.

Gli errori numerici possono quindi propagarsi lungo la traiettoria.

Diventano importanti:

\begin{itemize}
    \item la tolleranza del solver;
    \item il criterio di arresto;
    \item il numero massimo di iterazioni;
    \item la qualità del valore iniziale;
    \item la stabilità del problema dinamico.
\end{itemize}

Un sistema dinamicamente stabile tende, in alcune situazioni, ad attenuare
piccole perturbazioni numeriche; un sistema instabile può invece
amplificarle.

La stabilità economica e le proprietà numeriche dell'algoritmo possono
quindi interagire.


\subsection{Due esempi, due livelli di difficoltà}

Gli esempi precedenti mostrano due diversi casi di sistema implicito.

Nel modello SP-DG, la simultaneità può essere eliminata analiticamente:
\[
G(\mathbf{x}_{t+1},\mathbf{x}_t)=0
\quad\longrightarrow\quad
\mathbf{x}_{t+1}=F(\mathbf{x}_t).
\]

Il problema computazionale viene quindi risolto una volta, nella fase di
manipolazione del modello.

Nel modello finanziario, invece, la relazione non lineare rende utile
risolvere numericamente
\[
G(\mathbf{x}_{t+1},\mathbf{x}_t)=0
\]
a ogni periodo.

Possiamo quindi distinguere:
\[
\boxed{
\begin{array}{c}
\text{simultaneità eliminabile analiticamente}\\
\Downarrow\\
\text{ricorsione diretta}
\end{array}
}
\]
e
\[
\boxed{
\begin{array}{c}
\text{simultaneità da risolvere numericamente}\\
\Downarrow\\
\text{solver a ogni periodo}
\end{array}
}.
\]

Questo confronto mostra ancora una volta come modelli economicamente
differenti possano richiedere tecniche computazionali differenti pur
appartenendo alla stessa classe generale di sistemi dinamici impliciti.


\subsection{Dai sistemi impliciti ai sistemi forward-looking}

Finora, anche quando il sistema era implicito, dato lo stato corrente era
possibile determinare lo stato successivo risolvendo un problema
intra-periodale:
\[
G(\mathbf{x}_{t+1},\mathbf{x}_t)=0.
\]

Una difficoltà qualitativamente diversa emerge quando la determinazione
delle variabili correnti dipende dalle condizioni future del sistema.

In questo caso possono comparire relazioni del tipo
\[
G(E_t\mathbf{x}_{t+1},
\mathbf{x}_t,
\mathbf{x}_{t-1})=0.
\]

La difficoltà non consiste più soltanto nel risolvere un sistema di
equazioni al periodo corrente. Occorre determinare una traiettoria
intertemporale coerente con le condizioni forward-looking.

Nel caso dei modelli con aspettative razionali, questo conduce alla
distinzione tra variabili predeterminate e variabili forward-looking e,
nei sistemi lineari, a problemi basati sugli autovalori generalizzati.

Prima di affrontare questa ulteriore complicazione, completeremo però lo
studio dei sistemi deterministici considerando il caso non lineare e
mostrando come l'analisi locale della stabilità possa essere ricondotta,
ancora una volta, al calcolo degli autovalori.

\section{Modelli dinamici forward-looking}

Finora abbiamo considerato sistemi nei quali, dato lo stato corrente,
lo stato del periodo successivo poteva essere determinato direttamente
oppure risolvendo un sistema di equazioni intra-periodale.

Nel caso esplicito:
\[
\mathbf{x}_{t+1}=F(\mathbf{x}_t),
\]
lo stato successivo è ottenuto mediante una semplice valutazione della
legge di movimento.

Nel caso implicito:
\[
G(\mathbf{x}_{t+1},\mathbf{x}_t)=0,
\]
dato $\mathbf{x}_t$, è necessario risolvere un sistema di equazioni per
determinare $\mathbf{x}_{t+1}$.

In entrambi i casi, tuttavia, una volta noto lo stato corrente il
problema può essere affrontato procedendo periodo per periodo.

Una difficoltà differente emerge quando alcune variabili correnti
dipendono da valori futuri del sistema. Questi modelli vengono indicati
come \emph{forward-looking}.

Una rappresentazione schematica può essere
\[
G(\mathbf{x}_{t+1},\mathbf{x}_t,\mathbf{x}_{t-1})=0.
\]

La presenza di $\mathbf{x}_{t+1}$ in una condizione che determina
$\mathbf{x}_t$ modifica profondamente la natura del problema
computazionale. Il valore corrente non può essere determinato
utilizzando esclusivamente la storia passata del sistema.


\subsection{Perché la forward iteration non è sufficiente}

Consideriamo una semplice relazione
\[
q_t
=
x_t+\beta q_{t+1},
\]
dove $q_t$ è una variabile forward-looking e $x_t$ è una variabile
determinata dallo stato corrente.

Supponiamo inoltre che
\[
x_{t+1}=a x_t.
\]

Dato $x_0$, possiamo calcolare immediatamente
\[
x_1=ax_0,
\]
poi
\[
x_2=a x_1,
\]
e così via.

Per determinare $q_0$, invece, abbiamo bisogno di conoscere $q_1$:
\[
q_0=x_0+\beta q_1.
\]

Ma
\[
q_1=x_1+\beta q_2,
\]
e quindi il calcolo di $q_1$ richiede $q_2$.

Analogamente,
\[
q_2=x_2+\beta q_3.
\]

Otteniamo quindi una sequenza del tipo
\[
q_0
\longleftarrow
q_1
\longleftarrow
q_2
\longleftarrow
q_3
\longleftarrow
\cdots.
\]

La semplice forward iteration utilizzata nei sistemi ricorsivi non
chiude quindi il problema.


\subsection{Soluzione forward}

Possiamo sostituire ricorsivamente i valori futuri.

Partendo da
\[
q_t
=
x_t+\beta q_{t+1},
\]
utilizziamo
\[
q_{t+1}
=
x_{t+1}+\beta q_{t+2}
\]
e otteniamo
\[
q_t
=
x_t+\beta x_{t+1}
+\beta^2q_{t+2}.
\]

Continuando:
\[
q_t
=
x_t
+
\beta x_{t+1}
+
\beta^2x_{t+2}
+
\cdots
+
\beta^Tq_{t+T}.
\]

Se imponiamo una condizione che escluda una componente esplosiva della
soluzione, il termine terminale tende a zero e possiamo ottenere
\[
q_t
=
\sum_{j=0}^{\infty}
\beta^j x_{t+j}.
\]

La variabile corrente dipende quindi dall'intera sequenza futura della
variabile fondamentale.

Questo è il tratto caratteristico dei modelli forward-looking:
\[
\boxed{
\text{il presente dipende dalle condizioni future}
}.
\]


\subsection{Un caso con soluzione analitica}

Supponiamo che
\[
x_{t+1}=a x_t.
\]

Allora
\[
x_{t+j}=a^j x_t.
\]

La soluzione forward diventa
\[
q_t
=
x_t
\sum_{j=0}^{\infty}
(\beta a)^j.
\]

Se
\[
|\beta a|<1,
\]
la serie geometrica converge e otteniamo
\[
q_t
=
\frac{x_t}{1-\beta a}.
\]

Possiamo quindi scrivere
\[
q_t=P x_t,
\]
dove
\[
P=\frac{1}{1-\beta a}.
\]

Una volta trovata questa relazione, il modello può nuovamente essere
simulato procedendo in avanti.

Il problema computazionale è quindi composto da due fasi:
\[
\boxed{
\text{SOLVE}
\longrightarrow
\text{SIMULATE}
}.
\]


\subsection{Sistemi forward-looking e condizioni terminali}

Nei sistemi forward-looking la condizione iniziale non è generalmente
sufficiente a determinare la soluzione.

Nei modelli ricorsivi avevamo un Initial Value Problem:
\[
\begin{cases}
\mathbf{x}_{t+1}=F(\mathbf{x}_t),\\
\mathbf{x}_0=\bar{\mathbf{x}}.
\end{cases}
\]

Il passato determina lo stato iniziale e la legge di movimento permette
di procedere in avanti.

Nei sistemi forward-looking compaiono invece anche condizioni che
riguardano il comportamento futuro del sistema.

Queste possono assumere la forma di:

\begin{itemize}
    \item una condizione terminale;
    \item una condizione di non esplosività;
    \item una condizione di trasversalità;
    \item una condizione di equilibrio sulle aspettative.
\end{itemize}

Il problema assume quindi alcune caratteristiche tipiche di un
\emph{boundary value problem}: alcune condizioni provengono dal passato,
mentre altre limitano il comportamento futuro della traiettoria.


\subsection{Variabili predeterminate e variabili forward-looking}

Nei modelli intertemporali è utile distinguere tra due tipi di variabili.

Una variabile \emph{predeterminata} è determinata dalla storia del sistema.
Il suo valore al tempo $t$ è già noto quando vengono prese le decisioni
del periodo.

Un esempio tipico è uno stock:
\[
k_t.
\]

Il valore di $k_t$ dipende dalle decisioni e dagli eventi dei periodi
precedenti e non può generalmente modificarsi istantaneamente.

Una variabile \emph{forward-looking}, invece, può reagire immediatamente
alle informazioni sulle condizioni future.

Un prezzo finanziario costituisce un esempio naturale: nuove informazioni
sui dividendi o sui tassi futuri possono modificare immediatamente il
prezzo corrente.

La distinzione tra variabili predeterminate e forward-looking è centrale
per determinare quali traiettorie siano compatibili con l'equilibrio del
modello.


\subsection{Un semplice sistema lineare forward-looking}

Consideriamo, a titolo illustrativo, un sistema composto da una variabile
predeterminata $x_t$ e da una variabile forward-looking $q_t$:
\[
x_{t+1}=a x_t,
\]
\[
q_t=x_t+\beta q_{t+1}.
\]

Risolvendo la seconda relazione rispetto a $q_{t+1}$:
\[
q_{t+1}
=
\frac{1}{\beta}q_t
-
\frac{1}{\beta}x_t.
\]

Possiamo quindi scrivere
\[
\begin{bmatrix}
x_{t+1}\\
q_{t+1}
\end{bmatrix}
=
\begin{bmatrix}
a&0\\
-\frac{1}{\beta}&\frac{1}{\beta}
\end{bmatrix}
\begin{bmatrix}
x_t\\
q_t
\end{bmatrix}.
\]

Gli autovalori della matrice sono
\[
\lambda_1=a
\]
e
\[
\lambda_2=\frac{1}{\beta}.
\]

Se
\[
|a|<1
\]
e
\[
0<\beta<1,
\]
abbiamo
\[
|\lambda_1|<1
\]
ma
\[
|\lambda_2|>1.
\]

Se applicassimo la condizione utilizzata per i normali sistemi
ricorsivi, concluderemmo che il sistema è instabile.

Nei modelli forward-looking questa conclusione sarebbe però incompleta.


\subsection{Il ruolo delle radici instabili}

La presenza di un autovalore fuori dal cerchio unitario non implica
necessariamente che l'equilibrio forward-looking sia esplosivo.

Nel precedente esempio, $x_t$ è predeterminata, mentre $q_t$ può
aggiustarsi immediatamente.

Il valore iniziale $q_0$ non è quindi arbitrariamente assegnato.
Deve essere scelto in modo che la componente associata alla radice
esplosiva venga eliminata.

La soluzione
\[
q_t
=
\frac{x_t}{1-\beta a}
\]
individua precisamente la traiettoria che soddisfa questa condizione.

La variabile forward-looking ``salta'' quindi sul percorso compatibile
con la soluzione non esplosiva.


\subsection{Dalla stabilità alla determinazione dell'equilibrio}

Nei sistemi ricorsivi lineari la domanda principale era:
\[
\text{tutte le radici sono all'interno del cerchio unitario?}
\]

La condizione
\[
\rho(A)<1
\]
garantiva la stabilità asintotica.

Nei sistemi forward-looking la domanda cambia.

Alcune radici possono trovarsi fuori dal cerchio unitario, perché le
variabili forward-looking possono aggiustarsi in modo da eliminare le
componenti esplosive.

Diventa quindi importante confrontare:

\[
\text{numero delle radici instabili}
\]
con
\[
\text{numero delle variabili forward-looking}.
\]

Questa idea costituisce la base delle condizioni di
\emph{Blanchard--Kahn} utilizzate nei modelli lineari con aspettative
razionali.


\subsection{Esistenza, unicità e indeterminatezza}

In una formulazione standard, e sotto opportune condizioni di regolarità,
possiamo distinguere tre casi.

Se il numero delle radici instabili coincide con il numero delle variabili
forward-looking, il modello può possedere una soluzione non esplosiva
unica.

Se le radici instabili sono troppo poche rispetto alle variabili
forward-looking, possono esistere molte traiettorie compatibili con
l'equilibrio. Si parla di \emph{indeterminatezza}.

Se le radici instabili sono troppe, può non esistere una traiettoria
non esplosiva compatibile con le condizioni del modello.

Il problema degli autovalori assume quindi una nuova funzione.

Nei sistemi ricorsivi veniva utilizzato soprattutto per studiare la
stabilità.

Nei sistemi forward-looking viene utilizzato anche per analizzare:
\[
\boxed{
\text{esistenza}
\quad
\text{e}
\quad
\text{unicità della soluzione}
}.
\]


\subsection{Aspettative razionali}

Finora abbiamo formulato il problema in termini deterministici.

Quando viene introdotta l'incertezza, le variabili future non sono note con
certezza. Le condizioni forward-looking coinvolgono quindi aspettative
condizionate:
\[
E_t\mathbf{x}_{t+1}.
\]

Una possibile rappresentazione è
\[
A E_t\mathbf{x}_{t+1}
=
B\mathbf{x}_t
+
C\boldsymbol{\varepsilon}_t.
\]

Nei modelli con aspettative razionali, le aspettative utilizzate dagli
agenti devono essere coerenti con la distribuzione delle variabili
generata dal modello stesso.

Non è quindi possibile specificare arbitrariamente
\[
E_t\mathbf{x}_{t+1}.
\]

La soluzione deve essere tale che le aspettative formulate all'interno
del modello coincidano con quelle implicate dalla stessa soluzione.


\subsection{Un nuovo problema agli autovalori}

Nei sistemi lineari standard avevamo
\[
\mathbf{x}_{t+1}=A\mathbf{x}_t
\]
e studiavamo gli autovalori della singola matrice $A$.

Nei modelli forward-looking si incontrano frequentemente sistemi della
forma
\[
A E_t\mathbf{x}_{t+1}
=
B\mathbf{x}_t.
\]

La dinamica dipende quindi dalla coppia di matrici $(A,B)$.

Gli autovalori rilevanti soddisfano
\[
\det(B-\lambda A)=0.
\]

Si tratta di un \emph{generalized eigenvalue problem}.

Il passaggio
\[
\text{eigenvalues}
\longrightarrow
\text{generalized eigenvalues}
\]
riflette quindi una differenza nella struttura economica del modello,
non soltanto una maggiore complessità matematica.


\subsection{La decomposizione QZ}

Il problema agli autovalori generalizzati può essere affrontato
numericamente attraverso la \emph{generalized Schur decomposition},
comunemente associata all'algoritmo QZ.

L'idea generale è trasformare la coppia di matrici $(A,B)$ in una forma
numericamente più conveniente, preservando le informazioni necessarie
per determinare gli autovalori generalizzati.

Nei modelli macroeconomici lineari con aspettative razionali questo
permette, tra le altre cose, di:

\begin{itemize}
    \item distinguere radici stabili e instabili;
    \item ordinare le radici rispetto al cerchio unitario;
    \item verificare le condizioni di determinazione;
    \item ricavare la legge di movimento della soluzione.
\end{itemize}

Dal punto di vista computazionale la logica generale è quindi
\[
\boxed{
\text{modello strutturale}
\longrightarrow
\text{QZ / generalized eigenvalues}
\longrightarrow
\text{solution rule}
\longrightarrow
\text{simulazione}
}.
\]


\subsection{Un esempio economico: inflazione forward-looking}

Consideriamo una relazione molto semplice nella quale l'inflazione
corrente dipende dall'inflazione futura attesa e dall'output gap:
\[
\pi_t
=
\beta E_t\pi_{t+1}
+
\kappa y_t.
\]

Questa struttura è qualitativamente diversa dalla regola di aspettative
adattive considerata nel modello SP-DG:
\[
\pi_t^e
=
\lambda\pi_{t-1}
+
(1-\lambda)\pi_{t-1}^e.
\]

Nel caso delle aspettative adattive, l'aspettativa corrente è ottenuta
esclusivamente da variabili passate. Il modello può quindi essere
simulato ricorsivamente.

Nella relazione forward-looking:
\[
\pi_t
=
\beta E_t\pi_{t+1}
+
\kappa y_t,
\]
invece, per determinare $\pi_t$ dobbiamo conoscere le aspettative sulla
dinamica futura dell'inflazione.

La semplice forward iteration non è quindi sufficiente.

Se, ad esempio, la dinamica dell'output gap fosse
\[
y_{t+1}=a y_t,
\]
una soluzione potrebbe essere cercata nella forma
\[
\pi_t=P y_t.
\]

Sostituendo:
\[
P y_t
=
\beta P a y_t
+
\kappa y_t.
\]

Pertanto
\[
P
=
\beta aP+\kappa,
\]
da cui
\[
P
=
\frac{\kappa}{1-\beta a}.
\]

Una volta determinato $P$, il modello diventa simulabile:
\[
\pi_t
=
\frac{\kappa}{1-\beta a}y_t.
\]

Ancora una volta:
\[
\boxed{
\text{SOLVE}
\longrightarrow
\text{SIMULATE}
}.
\]


\subsection{Un esempio finanziario: il prezzo di un'attività}

I modelli di asset pricing forniscono un esempio particolarmente naturale
di struttura forward-looking.

In una formulazione estremamente semplice, il prezzo corrente di
un'attività può essere scritto come
\[
P_t
=
D_t+\beta E_tP_{t+1},
\]
dove $D_t$ rappresenta un payoff corrente e $\beta$ un fattore di sconto.

Iterando in avanti:
\[
P_t
=
D_t
+
\beta E_tD_{t+1}
+
\beta^2E_tP_{t+2}.
\]

Continuando:
\[
P_t
=
D_t
+
\beta E_tD_{t+1}
+
\beta^2E_tD_{t+2}
+
\cdots.
\]

Sotto una condizione di non esplosività appropriata, otteniamo
\[
P_t
=
\sum_{j=0}^{\infty}
\beta^j E_tD_{t+j}.
\]

Il prezzo corrente incorpora quindi informazioni relative all'intera
sequenza dei payoff futuri attesi.

Questo spiega perché i prezzi finanziari possono reagire immediatamente a
informazioni riguardanti eventi futuri.


\subsection{Una notizia futura può modificare il prezzo corrente}

Supponiamo, per semplicità, che tutti i payoff siano nulli ad eccezione
di un pagamento noto al periodo $T$:
\[
D_T=D>0.
\]

La soluzione è
\[
P_t
=
\beta^{T-t}D,
\qquad t\leq T.
\]

In particolare:
\[
P_0=\beta^T D.
\]

Il pagamento avverrà soltanto nel periodo $T$, ma il prezzo reagisce già
nel periodo corrente.

Questa è una conseguenza diretta della natura forward-looking del prezzo:
\[
\boxed{
\text{informazione sul futuro}
\longrightarrow
\text{aggiustamento immediato del presente}
}.
\]

Una normale legge ricorsiva backward-looking non presenterebbe
necessariamente questa caratteristica.


\subsection{Il significato computazionale del confronto}

Possiamo ora confrontare le tre strutture considerate finora.

Nel sistema ricorsivo:
\[
\mathbf{x}_{t+1}=F(\mathbf{x}_t),
\]
il computer procede direttamente da un periodo al successivo.

Nel sistema implicito:
\[
G(\mathbf{x}_{t+1},\mathbf{x}_t)=0,
\]
il computer deve risolvere un problema numerico a ogni periodo, ma la
transizione rimane determinabile localmente.

Nel sistema forward-looking:
\[
G(E_t\mathbf{x}_{t+1},\mathbf{x}_t,\ldots)=0,
\]
la soluzione corrente dipende dalla struttura futura del modello.

La differenza può quindi essere sintetizzata come
\[
\boxed{
\begin{array}{rcl}
\text{ricorsivo}
&:&
\text{iterate},\\[1mm]
\text{implicito}
&:&
\text{solve period by period},\\[1mm]
\text{forward-looking}
&:&
\text{solve the intertemporal model first}.
\end{array}
}
\]


\subsection{Dai modelli deterministici ai modelli stocastici}

I modelli forward-looking mostrano chiaramente perché la distinzione tra
deterministico e stocastico costituisca una dimensione separata della
classificazione.

In un modello deterministico con previsione perfetta, il futuro è
endogenamente determinato dal modello e gli agenti ne conoscono la
traiettoria.

In un modello stocastico, invece, il futuro non è noto. Gli agenti possono
soltanto formulare aspettative condizionate all'informazione corrente:
\[
E_t\mathbf{x}_{t+1}.
\]

L'introduzione dell'incertezza modifica quindi l'oggetto della soluzione:
non dobbiamo più determinare soltanto una traiettoria coerente, ma anche
specificare come gli shock generano una distribuzione delle possibili
traiettorie e come gli agenti formano aspettative su tale distribuzione.

Questo ci conduce naturalmente al passo successivo nello studio dei
modelli dinamici discreti: i modelli stocastici non controllati.

\section{Sintesi: struttura dinamica e problema computazionale}

L'analisi dei modelli deterministici in tempo discreto ha mostrato che
problemi economici e finanziari anche molto diversi possono essere
ricondotti a un numero relativamente limitato di strutture dinamiche.
La forma assunta dal modello determina, a sua volta, il tipo di problema
computazionale che deve essere affrontato.

Il caso più semplice è rappresentato dai sistemi ricorsivi espliciti:
\[
\mathbf{x}_{t+1}=F(\mathbf{x}_t).
\]
Una volta assegnata la condizione iniziale $\mathbf{x}_0$, lo stato
successivo può essere ottenuto direttamente valutando la funzione
$F$. La traiettoria viene quindi costruita mediante iterazione in avanti:
\[
\mathbf{x}_0
\rightarrow
\mathbf{x}_1
\rightarrow
\mathbf{x}_2
\rightarrow
\cdots.
\]
Dal punto di vista computazionale, il problema fondamentale è quindi
la simulazione della legge di movimento.

Nel caso lineare multidimensionale,
\[
\mathbf{x}_{t+1}=A\mathbf{x}_t+\mathbf{b},
\]
la stessa struttura ricorsiva permette di generare direttamente le
traiettorie, ma emerge un ulteriore problema: stabilire se tali
traiettorie convergano verso lo stato stazionario.

Definendo
\[
\widetilde{\mathbf{x}}_t
=
\mathbf{x}_t-\mathbf{x}^*,
\]
la dinamica delle deviazioni dallo steady state è
\[
\widetilde{\mathbf{x}}_{t+1}
=
A\widetilde{\mathbf{x}}_t,
\]
e quindi
\[
\widetilde{\mathbf{x}}_t
=
A^t\widetilde{\mathbf{x}}_0.
\]
La stabilità dipende pertanto dagli autovalori della matrice di
transizione. In particolare, nel caso lineare discreto lo steady state
è asintoticamente stabile se
\[
\rho(A)<1,
\]
dove $\rho(A)$ è il raggio spettrale della matrice.

L'analisi della stabilità introduce quindi un secondo importante
problema computazionale: il calcolo degli autovalori. Per matrici dense
di dimensione moderata possono essere utilizzati algoritmi basati sulla
decomposizione QR. Nei sistemi di grandi dimensioni, soprattutto quando
la matrice di transizione è sparsa, diventa invece conveniente
utilizzare metodi iterativi che sfruttano la struttura del problema,
come il power method e i metodi basati sui sottospazi di Krylov.

Una struttura diversa emerge nei sistemi dinamici impliciti:
\[
G(\mathbf{x}_{t+1},\mathbf{x}_t)=0.
\]
In questo caso, anche conoscendo $\mathbf{x}_t$, lo stato successivo non
è necessariamente disponibile mediante la semplice valutazione di una
funzione. Esso deve essere determinato risolvendo un sistema di
equazioni.

Nel caso lineare,
\[
A\mathbf{x}_{t+1}
=
B\mathbf{x}_t+\mathbf{c},
\]
la transizione richiede la soluzione di un sistema lineare. Nel caso
non lineare,
\[
G(\mathbf{x}_{t+1},\mathbf{x}_t)=0,
\]
essa può invece richiedere un algoritmo di ricerca degli zeri.

La struttura temporale rimane tuttavia quella di un initial value
problem, purché, dato $\mathbf{x}_t$, esista una soluzione unica per
$\mathbf{x}_{t+1}$. La differenza rispetto al sistema esplicito riguarda
quindi il modo in cui viene calcolata ciascuna transizione:
\[
\text{esplicito}
\quad\Rightarrow\quad
\text{evaluate},
\]
mentre
\[
\text{implicito}
\quad\Rightarrow\quad
\text{solve}.
\]

La situazione cambia più profondamente nei modelli forward-looking.
In questi modelli alcune variabili correnti dipendono dal comportamento
futuro del sistema. La sola condizione iniziale non è quindi, in
generale, sufficiente per determinare la traiettoria.

Occorre distinguere tra variabili predeterminate, il cui valore iniziale
è ereditato dalla storia del sistema, e variabili forward-looking, che
possono aggiustarsi immediatamente in modo da rendere la traiettoria
compatibile con le condizioni future.

Il problema assume quindi la struttura di un boundary value problem o,
più precisamente in molti modelli economici, di un mixed boundary value
problem. La soluzione deve soddisfare simultaneamente condizioni
iniziali e condizioni future, terminali o di non esplosività.

Dal punto di vista computazionale cambia pertanto la sequenza delle
operazioni. Nei sistemi ricorsivi possiamo scrivere schematicamente
\[
\text{SIMULATE},
\]
mentre nei sistemi forward-looking la sequenza diventa
\[
\text{SOLVE}
\rightarrow
\text{SIMULATE}.
\]
Prima di generare una traiettoria occorre cioè individuare la regola di
soluzione compatibile con le condizioni intertemporali del modello.

Anche il ruolo degli autovalori cambia. Nei sistemi ricorsivi essi
servono principalmente a stabilire se le perturbazioni tendano a
scomparire o ad amplificarsi. Nei modelli forward-looking, invece, la
separazione tra radici stabili e instabili contribuisce anche a
determinare l'esistenza e l'unicità della soluzione.

Nei sistemi lineari con aspettative razionali questo conduce alle
condizioni di Blanchard--Kahn e, più in generale, a problemi agli
autovalori generalizzati. Dal punto di vista numerico, strumenti come
la decomposizione QZ permettono di separare le radici stabili da quelle
instabili e di ricavare la regola di soluzione del modello.

Possiamo quindi sintetizzare il percorso seguito nel modo seguente:
\[
\begin{array}{ccl}
\text{sistema ricorsivo} & \longrightarrow &
\text{iterazione in avanti},\\
\text{sistema ricorsivo multidimensionale} & \longrightarrow &
\text{autovalori e stabilità},\\
\text{sistema implicito lineare} & \longrightarrow &
\text{soluzione di sistemi lineari},\\
\text{sistema implicito non lineare} & \longrightarrow &
\text{root finding},\\
\text{sistema forward-looking} & \longrightarrow &
\text{BVP e soluzione intertemporale},\\
\text{sistema con aspettative razionali} & \longrightarrow &
\text{autovalori generalizzati e QZ}.
\end{array}
\]

\subsection{Una precisazione sulla stabilità locale}

L'analisi svolta per i sistemi lineari ha un'importante proprietà. Se
\[
\mathbf{x}_{t+1}=A\mathbf{x}_t+\mathbf{b}
\]
e
\[
\rho(A)<1,
\]
la condizione sugli autovalori caratterizza direttamente la dinamica
globale del sistema lineare: qualunque deviazione iniziale dallo steady
state viene progressivamente riassorbita.

Nei sistemi non lineari la situazione è più complessa. Consideriamo
\[
\mathbf{x}_{t+1}=F(\mathbf{x}_t)
\]
e uno steady state $\mathbf{x}^*$ tale che
\[
\mathbf{x}^*=F(\mathbf{x}^*).
\]
Nelle vicinanze dello steady state possiamo approssimare il sistema
mediante la linearizzazione
\[
\mathbf{x}_{t+1}-\mathbf{x}^*
\simeq
J(\mathbf{x}^*)
(\mathbf{x}_t-\mathbf{x}^*),
\]
dove
\[
J(\mathbf{x}^*)
=
\left.
\frac{\partial F(\mathbf{x})}
{\partial\mathbf{x}}
\right|_{\mathbf{x}=\mathbf{x}^*}
\]
è la matrice Jacobiana valutata nello steady state.

Gli autovalori della Jacobiana permettono quindi di studiare il
comportamento delle traiettorie nelle vicinanze dell'equilibrio. Se
\[
\rho\!\left(J(\mathbf{x}^*)\right)<1,
\]
lo steady state è localmente asintoticamente stabile.

La parola \emph{localmente} è essenziale. La linearizzazione descrive
il comportamento del sistema nelle vicinanze dello steady state, ma
non necessariamente quello di traiettorie che partono molto lontano
da esso. Un sistema non lineare può infatti presentare più steady state,
regioni di attrazione differenti o dinamiche che non possono essere
dedotte dalla sola analisi locale.

Questo chiarisce anche la complementarità tra analisi della stabilità
e simulazione numerica. Gli autovalori della Jacobiana forniscono
informazioni rigorose sul comportamento locale, mentre la simulazione
permette di esplorare la dinamica del sistema anche lontano dallo
steady state.

\subsection{Dal deterministico allo stocastico}

Tutti i modelli considerati finora sono stati analizzati in un ambiente
deterministico. Una volta assegnati i parametri e le condizioni
necessarie per determinare la soluzione, il modello genera una
traiettoria:
\[
\mathbf{x}_0,
\mathbf{x}_1,
\mathbf{x}_2,
\ldots.
\]

L'introduzione dell'incertezza modifica questa struttura. Una legge di
movimento può assumere, ad esempio, la forma
\[
\mathbf{x}_{t+1}
=
F(\mathbf{x}_t,\boldsymbol{\varepsilon}_{t+1}),
\]
dove $\boldsymbol{\varepsilon}_{t+1}$ rappresenta uno shock casuale.

A parità di condizione iniziale e di parametri, realizzazioni differenti
degli shock generano traiettorie differenti:
\[
\mathbf{x}_0
\rightarrow
\mathbf{x}_1^{(1)}
\rightarrow
\mathbf{x}_2^{(1)}
\rightarrow\cdots,
\]
\[
\mathbf{x}_0
\rightarrow
\mathbf{x}_1^{(2)}
\rightarrow
\mathbf{x}_2^{(2)}
\rightarrow\cdots.
\]

L'oggetto di interesse non è quindi più soltanto una singola traiettoria
deterministica. Diventano rilevanti la distribuzione delle possibili
traiettorie, i loro momenti e la probabilità che il sistema raggiunga
determinate regioni dello spazio degli stati.

Anche il problema computazionale cambia. Alla semplice iterazione della
legge di movimento si aggiunge la necessità di generare realizzazioni
degli shock e di ripetere la simulazione per molte sequenze casuali.
Questo conduce naturalmente ai metodi di simulazione stocastica e,
in particolare, alle tecniche Monte Carlo.

Nei modelli forward-looking l'incertezza introduce inoltre un ulteriore
elemento. Le variabili correnti possono dipendere dalle aspettative
condizionate sui valori futuri:
\[
E_t\mathbf{x}_{t+1}.
\]
Nel caso delle aspettative razionali tali aspettative devono essere
coerenti con la distribuzione futura generata dal modello stesso.

Il passaggio dal deterministico allo stocastico non elimina quindi i
problemi computazionali incontrati finora. Al contrario, li combina con
nuovi problemi:
\[
\text{dinamica}
+
\text{incertezza}
+
\text{aspettative}.
\]

Questa sarà la struttura di riferimento per lo studio dei modelli
dinamici discreti stocastici.


\chapter{modelli stocastici}

\end{document}
